\section{Summary for stakeholders}
\label{sec:stakeholders}
% ============================================================================

\epigraph{``Security is a process, not a product.''}{--- Bruce Schneier, \textit{Secrets and Lies}~\cite{schneier2000}}

This report describes a security architecture model based on
\emph{segmentation, Trusted Activity Contexts, and federated
governance}, applicable to research-and-teaching universities under
the \gls{nis2} essential-entity regime. The messages below are
tailored to each category of stakeholders to facilitate internal
communication.

\begin{tcolorbox}[colback=Slate!4!white, colframe=Slate!60!black,
  sharp corners=south, boxrule=0.5pt, breakable,
  left=6pt, right=6pt, top=5pt, bottom=5pt]
\textit{\textbf{On the nature of this report.} The present
document is a voluntary contribution produced by members of the
institution's academic community acting in
their scientific and technical capacity. It is transversal across
the institutional actors concerned (\loc{executive-leadership}, faculties,
interdisciplinary centres, services, IT departments, DPO, CISO)
and independent of any institutional hierarchy or formal mandate.
The single objective of this report is to be considered by the
institution's academic governing bodies (\loc{academic-governing-bodies} and
\loc{governing-board}) and, if found substantive, to
be propagated by \loc{governing-board} to all concerned components
for deliberation through the appropriate institutional channels.}
\end{tcolorbox}

% --- Top governing board / executive leadership ---
\begin{stakeholderbox}[Message to \loc{governing-board} and \loc{executive-leadership}]{blue}
\textbf{Strategic issue:} The university's capacity to attract
and retain world-class researchers, to compete for European and
international funding, and to meet its legal obligations under
EU digital regulation --- all three are directly shaped by its
digital infrastructure choices. In accordance with the
\emph{Magna Charta Universitatum} and the UNESCO \emph{Recommendation
on Science and Scientific Researchers} (2017), digital openness is not
a concession but a founding condition of the university.
A university requires administration, but is not a corporation;
university governance must optimise for freedom for ideas and persons,
physical or digital. The digital infrastructure should reflect this
fundamental distinction.

\medskip
Article~32 of the \gls{gdpr} requires security measures \emph{appropriate to
the risk}, including \keyword{resilience} and the capacity for
\keyword{restoration}~\cite{gdpr}. Where the national transposition of
the \gls{nis2} Directive~\cite{nis2} so provides, an institution of this kind is classified
by the national competent authority --- the NIS2 supervisory authority --- as an
\keyword{essential entity} on the basis of its status as a public
institution; this classification applies to the institution as a whole
and entails proactive supervision, regular audits and an annual
compliance and maturity report. Sector guidance
treats cybersecurity as a \emph{strategic governance risk} and
recommends defence in depth and a positive security
culture~\cite{jisc-guidance}.

\medskip
The analysis of the most publicised university breaches in
2024--2025 (Paris-Saclay, Texas Tech, Western Sydney ---
see Section~\ref{sec:defaillances-mediatiques}) reveals that
\emph{none} were caused by research-zone autonomy or segmented
architectures; all stemmed from centralised, unsegmented
perimeters, supply-chain vulnerabilities, or credential
compromise.

\medskip
\textbf{What we propose:}
\begin{itemize}[nosep]
    \item \textbf{Multi-zone architecture with Trusted Activity
    Contexts:} a segmented infrastructure that \emph{strengthens}
    the protection of sensitive administrative data by isolating
    them in a dedicated enclave (recent media-covered incidents
    show that centralised, undifferentiated perimeters amplify
    impact --- see Section~\ref{sec:defaillances-mediatiques}),
    while each research project, teaching activity, or
    administrative process operates within a validated, isolated
    environment adapted to its risk level --- ensuring that an
    incident in one context cannot compromise the rest of the
    institution;
    \item \textbf{Federated governance with responsibility charters}
    aligned with the models of world-class universities (MIT,
    Oxford, Cambridge, ETH Zurich, Michigan), formalising autonomy
    and accountability (ETH~Zurich model) and replacing informal
    exceptions;
    \item \textbf{Auditable regulatory compliance:} a programme
    producing evidence of effectiveness compliant with the
    \gls{gdpr}, \gls{nis2} and the European AI Act, preserving
    attractiveness for high-level researchers and academics
    through a modern infrastructure;
    \item \textbf{AI as a structural research tool (2025--2026):}
    ensuring access to major cloud-based AI services (OpenAI,
    Anthropic, Google, HuggingFace, AWS) --- under the AI Act's
    risk-based framework, the applicable obligations depend on
    the concrete use case and the user's role; many research
    uses of commercial AI services fall outside the high-risk
    categories;
    \item An \textbf{Independent Commission for Reasonable
    Adjustments of Digital Policies} (\gls{icradp}), guaranteeing
    every employee the right to obtain adjustments for medical,
    personal, professional reasons or infringement of fundamental
    freedoms --- in accordance with the EU Charter of Fundamental
    Rights, the CRPD Convention and the case law of the \gls{cjeu}
    and the \gls{ecthr} on proportionality (where these instruments
    bind the institution's jurisdiction)
    (see Section~\ref{sec:commission-ajustements}).
\end{itemize}

\medskip
\textbf{Request for consideration by the academic governing bodies and
propagation by \loc{governing-board}.} The present
report is a voluntary contribution produced by members of the
academic community acting in their scientific and technical
capacity, independent of any institutional hierarchy or formal
mandate. Its propositions are offered to the institution's academic
governing bodies (\loc{academic-governing-bodies} and \loc{governing-board}) for
consideration alongside the cybersecurity strengthening initiatives
already approved by \loc{governing-board}. The contributors of
this report request that, if the analysis is found substantive,
the document be propagated by \loc{governing-board} to all
concerned components (\loc{executive-leadership}, faculties, interdisciplinary
centres, services, IT departments, DPO, CISO) so that integration
with existing initiatives can be deliberated through the
appropriate institutional channels.

\medskip
\textbf{Feasibility:} the proposed measures can be implemented
\emph{progressively} (governance first, then pilot, then rollout)
within a university multi-year plan, leveraging existing
\gls{nren}/\gls{geant} services, open-source tools and shared
security communities (G\'EANT Security, the national NREN's CSIRT,
\gls{first}). A digital infrastructure worthy of the university's
ambitions is a strategic enabler, not merely a cost
(see Section~\ref{sec:couts} and Section~\ref{sec:transition}).

\medskip
\textbf{Documented risk factors:} talent mobility patterns toward
more flexible institutions; regulatory exposure under \gls{nis2}
essential-entity supervision; Shadow IT as a documented effect of
restrictive policies (a G\'EANT survey reports that in most
institutions \emph{``all official IT can have a shadow
version''}~\cite{geant-shadow-it}; Gartner projects that 75\% of
employees will acquire technology outside IT visibility by
2027~\cite{gartner-shadow-it}); competitiveness factors in
European and international grant calls; potential litigation based
on fundamental rights (academic freedom, freedom of expression and
information, dignity at work). The analysis of the most publicised
university failures in 2024--2025
(Section~\ref{sec:defaillances-mediatiques}) identifies third-party
vulnerabilities, credential compromise and absence of internal
boundaries as recurring causes.

\medskip
\textbf{Proposed next steps:}
\begin{enumerate}[nosep]
    \item The report is submitted to the institution's academic
    governing bodies (\loc{academic-governing-bodies} and
    \loc{governing-board}) for consideration within
    their respective remits; each body may issue a position
    statement on the propositions.
    \item If endorsed by \loc{governing-board}, \loc{executive-leadership}
    --- in collaboration with the CISO, the DPO, the academic
    representatives and the contributors of this report ---
    defines an integrated implementation roadmap consistent with
    the board-approved initiatives (governance charter first, then
    pilot on selected entities, then progressive rollout).
    \item The contributors remain available, post-deliberation,
    to support the mapping between the propositions endorsed by
    the academic governing bodies and the integrated implementation
    roadmap.
\end{enumerate}
See the accompanying \emph{Executive Summary} for the detailed
call to action.
\end{stakeholderbox}

\bigskip

% --- CISO ---
\begin{stakeholderbox}[Message to the CISO]{red}
\textbf{Issue:} Reducing the attack surface you must defend,
gaining visibility on what is currently Shadow IT, and
obtaining auditable proof that security controls are effective
--- not just deployed.

\medskip
The core proposal is a \textbf{net security gain}: the
administrative enclave becomes \emph{smaller, easier to harden
and easier to monitor}. NIST SP~800-53 SC-7
(\emph{Boundary Protection})~\cite{nist-sp800-53} explicitly
requires organisations to ``connect to external networks or
information systems only through managed interfaces consisting
of boundary protection devices''; SC-7(21) calls for
\emph{isolation of system components}. ISO~27001 Annex~A
control A.8.22 requires \emph{network
segregation}~\cite{iso27001}. ETH Zurich~\cite{eth-baseline}
requires that non-compliant IT resources be placed in
\emph{isolated network zones}. High-profile university incidents in 2024--2025
(analysed in Section~\ref{sec:defaillances-mediatiques}) confirm that
breaches are driven by \emph{lack} of segmentation and lateral
movement, not by accountable autonomy in research or teaching zones.

\medskip
\textbf{What changes for you:}
\begin{itemize}[nosep]
    \item Research/teaching autonomy is governed by
    \emph{charters}, \emph{inventories} and \emph{monitoring} ---
    not by informal exceptions;
    \item Each zone has an identified owner, incident contacts,
    and a logging baseline;
    \item Each activity operates within a \emph{Trusted Activity
    Context} with its own security policies, logging baseline, and
    isolation boundaries --- making audit and incident response more
    precise and efficient;
    \item \textbf{Measurable security posture:} each Trusted
    Activity Context produces its own logging baseline and
    compliance metrics --- enabling you to report zone-by-zone
    security status to \loc{governing-board}, rather than a
    single aggregate score that masks local weaknesses;
    \item Current Shadow IT (hidden servers, personal clouds, unauthorised
    \gls{vpn}s) is replaced by authorised and auditable pathways;
    \item Oxford~\cite{oxford-infosec} and ETH Zurich~\cite{eth-baseline}
    provide credible precedents for this accountable autonomy.
\end{itemize}

\medskip
\textbf{Risk of other models for the CISO:} a single,
undifferentiated perimeter means that every exception is ad hoc,
every research server is a potential blind spot, and a breach
anywhere is a breach everywhere. The Verizon
DBIR~\cite{verizon-dbir2024} consistently identifies credential
abuse and ransomware as leading breach patterns; NIST guidance
identifies segmentation as the control that interrupts their
lateral progression across a flat network~\cite{nist-sp800-215}.
The proposed architecture does not weaken your mandate ---
it makes it enforceable.
\end{stakeholderbox}

\bigskip

% --- Central IT Service ---
\begin{stakeholderbox}[Message to the Central IT Department (CIO Office)]{teal}
\textbf{Issue:} Reducing the perimeter to manage, clarifying
responsibilities, and ending ad hoc exceptions.

\medskip
Today, a central IT Department that attempts to manage a single
perimeter covering the entire institution --- from financial data to
GPU compute clusters creates an unsustainable workload and
permanent friction with academic departments.
The consequence is a team stretched across incompatible
requirements: hardening administrative systems while
simultaneously troubleshooting research instruments,
configuring teaching labs, and managing ad hoc exceptions
that bypass the very policies the team is tasked to enforce.
Recent media-covered
incidents (e.g.\ Paris-Saclay, Western Sydney) show that such
undifferentiated perimeters also \emph{increase impact} when a breach
occurs --- segmentation limits propagation (Section~\ref{sec:defaillances-mediatiques}).

\medskip
\textbf{What changes for you:}
\begin{itemize}[nosep]
    \item You focus on the \emph{administrative enclave} (ERP,
    HR, finance, student IS) with strict controls and a well-defined
    perimeter;
    \item You provide the \emph{backbone}, \emph{federated identity}
    (\gls{iam}/\gls{mfa}) and the \emph{\gls{soc}} as shared services;
    \item Each department's research and teaching activities
    operate within \emph{Trusted Activity Contexts} --- you
    provision the underlying infrastructure (\gls{vlan}, identity,
    monitoring), but the domain-specific configuration is the
    department's responsibility under charter;
    \item Departments manage their own research servers in
    dedicated \gls{vlan}s, under their responsibility formalised by charter;
    \item The \emph{service catalogue} clarifies who provides what,
    replacing case-by-case negotiations;
    \item \textbf{No more ad hoc exceptions:} what is currently
    handled through informal email threads and individual
    workarounds becomes a formal charter process with documented
    responsibilities --- reducing your liability and your team's
    interrupt-driven workload;
    \item The University of Wisconsin-Madison~\cite{uw-madison} confirms
    that federation \emph{reduces risk} while preserving
    local innovation.
\end{itemize}

\medskip
\textbf{What does not change:} you retain full authority over
the backbone, the identity infrastructure, the SOC, and the
administrative enclave. The proposal does not reduce your
mandate --- it focuses it on the systems where central control
is both necessary and effective.
\end{stakeholderbox}

\bigskip

% --- Research Units and Research-and-Teaching Departments ---
\begin{stakeholderbox}[Message to Research Units and Research-and-Teaching Departments]{OliveGreen}
\textbf{Issue:} Regaining the autonomy necessary for research and
experimentation, whether in research-only units or research-and-teaching
departments. Concerns that ``more autonomy means more risk'' are
not supported by the evidence: the most publicised university breaches
in 2024--2025 (MOVEit, Oracle EBS, credential compromise at central
services) stem from \emph{third-party vulnerabilities} and
\emph{centralised perimeters}, not from research zones (see
Section~\ref{sec:defaillances-mediatiques}).

\medskip
\textbf{What changes for you:}
\begin{itemize}[nosep]
    \item Each laboratory or department can obtain a
    \emph{Trusted Activity Context} --- a dedicated \gls{vlan} with local
    administration rights (software installation, server configuration,
    root access) under a responsibility charter;
    \item Each unit designates a \emph{local IT contact}
    responsible for maintaining the charter commitments
    (inventory, patching, incident reporting) --- this role
    is the counterpart of the autonomy granted;
    \item Access to global cloud services (including for AI:
    OpenAI API, Google Cloud AI, AWS SageMaker, HuggingFace) is possible
    from the research zone, with an appropriate data management framework;
    \item Research instruments, laboratory IoT and GPU clusters
    are managed locally; high-throughput research traffic benefits
    from a \emph{Science DMZ} path bypassing the enterprise
    firewall (see Section~\ref{subsec:science-dmz});
    \item The model draws on ETH Zurich where an \emph{administrator rights
    agreement}~\cite{eth-admin-agreement} formalises root access
    for researchers and doctoral students;
    \item Research platforms such as CloudLab~\cite{cloudlab}
    and Chameleon~\cite{chameleon} illustrate that privileged
    control is a \emph{structural need} of computer science
    research;
    \item If a specific digital policy still restricts a
    legitimate research need, the \emph{Independent Commission}
    (\gls{icradp}) guarantees a formal right to request
    adjustments (see Section~\ref{sec:commission-ajustements}).
\end{itemize}
\end{stakeholderbox}

\bigskip

% --- Teaching and Training Components ---
\begin{stakeholderbox}[Message to Teaching and Training Components]{orange}
\textbf{Issue:} Providing teaching environments adapted to the
actual needs of each discipline --- from configurable labs for
computer science and engineering, to AI-assisted tools for law,
humanities and social sciences, to simulation platforms for
health sciences and economics. In 2025--2026, every faculty
needs access to cloud-based AI services and discipline-specific
digital tools; a uniform, locked-down desktop policy cannot
serve this diversity.

\medskip
At world-class universities with federated IT governance (MIT,
Stanford, ETH Zurich, Michigan), each school or department manages
its own IT infrastructure --- including the environments used for
teaching (see Section~\ref{sec:etat-art}). The teaching zone
proposed here applies this principle systematically.

\medskip
\textbf{What changes for you:}
The teaching zone provides \emph{Trusted Activity Contexts} ---
isolated, pre-configured environments adapted to each course or
activity:
\begin{itemize}[nosep]
    \item \gls{vdi} (\emph{Virtual Desktop Infrastructure}) allows
    students across all faculties to access standardised,
    discipline-specific environments from their
    own devices (\gls{byod}); students on personal devices access
    the same environments as in-lab students, ensuring equity
    between on-campus and remote learners;
    \item Access to AI APIs (ChatGPT, Claude, Gemini, open-source
    models) is available for all disciplines --- from
    AI-assisted legal research to computational linguistics,
    from data analysis in social sciences to creative tools
    in humanities;
    \item JupyterHub on Kubernetes~\cite{z2jh} provides each student
    with an isolated computing environment for data science,
    statistics, and computational courses across faculties;
    \item \emph{Cyber-ranges} (isolated sandbox environments) are
    made available for security, networking and systems courses,
    as recommended by NIST~\cite{nist-cyberrange};
    \item Instructors in every faculty can freely configure
    lab environments in the teaching zone, with a fast
    \emph{reset/reimage} pipeline between sessions;
    \item Instructors can prepare and test course environments
    \emph{before} the semester begins, with the certainty that
    the configuration will be preserved --- no more last-minute
    requests to central IT for software installation.
\end{itemize}
\end{stakeholderbox}

\bigskip

% --- Staff and BYOD ---
\begin{stakeholderbox}[Message to all Staff (BYOD)]{purple}
\textbf{Issue:} Being able to use one's own devices in a productive
and secure manner. The analysis of recent university incidents
(Section~\ref{sec:defaillances-mediatiques}) does not attribute breaches
to \gls{byod} or flexible access but to supply-chain, credentials and
unsegmented perimeters --- the proposed \gls{byod} zone with progressive
interoperability is designed to integrate personal devices safely
rather than push them into shadow IT.

\medskip
\textbf{What changes for you:}
\begin{itemize}[nosep]
    \item \textbf{Your personal device works on campus:} a
    dedicated network provides full Internet access from day one;
    as you authenticate, you progressively gain access to
    university resources matching your role (teaching materials,
    research data, administrative tools);
    \item With \gls{mfa} authentication and minimal posture verification
    (up-to-date OS, active antivirus), you can access
    institutional resources from your personal device;
    \item Access to sensitive data (HR, finance) from a \gls{byod}
    device remains conditional on enhanced controls (\gls{vdi}, secure
    web access);
    \item Access to cloud-based AI services (ChatGPT, Claude,
    Gemini) is available from the \gls{byod} network --- you do
    not need a university-managed device to use AI tools for
    your work;
    \item \gls{eduroam}~\cite{eduroam} already provides the federated
    Wi-Fi infrastructure: your local credentials work at all
    partner universities worldwide;
    \item \textbf{Right to reasonable adjustments:} an
    \emph{Independent Commission} (\gls{icradp}) guarantees you
    the right to request adjustments to any digital policy ---
    for example, if a software restriction prevents you from
    using an assistive technology, or if a network filter blocks
    a resource essential to your research. The commission is
    independent of the IT Department and rules within
    \loc{icradp-decision-deadline}
    (see Section~\ref{sec:commission-ajustements}).
\end{itemize}
\end{stakeholderbox}

% ============================================================================
\section{Reading guide: regulation, academic mission, and the institutional context}
\label{sec:reading-guide}
% ============================================================================

This section provides a pedagogical overview intended for
non-specialist decision-makers, particularly members of \loc{academic-governing-bodies} and
\loc{governing-board}. It introduces the regulatory
landscape, the academic-domain specificities, the institutional
context and the architectural proposition of this report. Each
sub-section is designed for a five-minute read; cross-references
point to the detailed treatment elsewhere in the report. Readers
familiar with the regulatory framework may proceed directly to
Section~\ref{sec:reglementation}.

\subsection{What this report is about}

This report is a voluntary contribution from members of the
institution's academic community, addressed to
the academic governing bodies for consideration (see ``On the nature of
this report'' at the entry of
Section~\ref{sec:stakeholders}). It describes a security
architecture model for the institution's information systems
--- segmentation, Trusted Activity Contexts and federated
governance --- and the operational mechanisms (researcher-centred
AI cloud access, institutional maturity reporting, executive
training) that align the architecture with the \gls{nis2}
essential-entity regime where the institution's jurisdiction has
transposed the NIS2 Directive into national law.

\subsection{The academic domain in five minutes}

A research-and-teaching university operates under three concurrent
missions, each with distinct working requirements:

\begin{itemize}[nosep]
    \item \textbf{Research.} Experimentation, autonomy, fast
    iteration; tools chosen by the researcher; international
    collaboration; access to high-performance computing and global
    data resources.
    \item \textbf{Teaching.} Structured curricula; access to
    industry-standard tools (including AI services such as
    OpenAI, Anthropic and Google AI for hands-on exercises);
    reproducible learning environments.
    \item \textbf{Administration.} Structured processes; strict
    protection of financial, HR and student records;
    regulated workflows.
\end{itemize}

\noindent The architectural challenge is that these three
missions have different requirements: research benefits from
openness and autonomy, administration requires strict control. A
single uniform policy applied to all three concentrates risk on a
single perimeter and constrains research workflows. The
multi-zone architecture of this report allocates a distinct
security profile to each mission, with explicit boundaries and
federated governance.

The institution does not operate in isolation: its national
research-and-education network (NREN)
connects it to its continental research backbone (in Europe,
\gls{geant}; Internet2/ESnet, RedCLARA, UbuntuNet or APAN in other
regions),
providing identical infrastructure primitives (eduroam, eduGAIN,
identity federations) at the national scale. The architectural
choices proposed in this report are consistent with the practice
of peer institutions accessible via the same backbone (see
Section~\ref{sec:exemples}).

\subsection{The European regulatory landscape}

For an adopter bound by the EU/EEA regime, four European regulatory
instruments shape the cybersecurity and data-protection environment
of the institution (an adopter under another jurisdiction
substitutes its own equivalents through the Regulatory Applicability
module of the Technical Reference):

\begin{description}[nosep]
    \item[\gls{gdpr} (2016, applicable since 2018).]
    Personal-data protection. \textbf{Article~32:} security
    measures \emph{appropriate to the risk}, including
    resilience and capacity for restoration.
    \textbf{Articles~24(1) and~5(2):} accountability principle
    --- the data controller is responsible. For a university,
    this applies to student records, HR data, and research data
    containing personal information.

    \item[\gls{nis2} (2022, transposition deadline 17~October~2024; national measures apply from 18~October~2024).]
    Cybersecurity for \emph{essential} and \emph{important}
    entities. Where the institution's jurisdiction has transposed
    the Directive into national law, a public research-and-teaching
    institution \loc{essential-entity-modality} as an \emph{essential entity}\loc{nis2-status-conditional-tail} --- the public-administration
    route being a common basis (Section~\ref{sec:nis2}).
    Essential-entity supervision is proactive (audits, maturity
    reports), with sanctions \loc{nis2-sanction-ceiling}.

    \item[European AI Act (Regulation~(EU)~2024/1689).]
    Risk-based classification of AI systems.
    \textbf{Article~2(6):} bespoke research-only AI tools are
    exempt. \textbf{Recital~25:} research and innovation
    activities are to be supported and not unduly restricted.
    For a university using commercial AI APIs (e.g.\ OpenAI,
    Anthropic, Google) in research and teaching, most uses fall
    outside the high-risk categories.

    \item[CJEU \emph{Schrems~II} (case C-311/18, 2020).]
    Transfers of personal data outside the European Economic Area
    require, in addition to Standard Contractual Clauses, a
    documented \gls{tia} and supplementary measures (encryption
    before transfer, pseudonymisation). This affects access to
    non-EU AI cloud services in research contexts (see
    Section~\ref{sec:safeguards-ai}).
\end{description}

\noindent\textit{Operational consequence.} For a university
classified as essential entity, the four instruments combine
into a single integrated regime: \gls{gdpr} sets the data-protection
baseline; NIS2 and its transposition set the cybersecurity
governance and maturity expectations; the AI Act sets the
AI-system classification; \emph{Schrems~II} constrains
cross-border data flows. \textcolor{Crimson}{\textbf{Where the adopter is bound by these
instruments, the architecture of this report is designed to satisfy
them simultaneously; an adopter under another regime reads the same
architecture against its own equivalents.}}

\subsection{The institution-specific context}

The elements below describe how the EU-level regime is typically
instantiated at national level. They are stated in role-based terms;
each institution substitutes the named authority, platform and
statute applicable in its own jurisdiction (the worked instantiation
is given in the Regulatory Applicability annex).

\begin{description}[nosep]
    \item[Status.] A public
    research-and-teaching institution \loc{essential-entity-modality} as an
    \emph{essential entity}\loc{nis2-status-conditional-tail}; the
    classification rests on three cumulative legal pegs:
    the public-administration entity definition, the essential-%
    regardless-of-size clause, and the research-organisation track
    (which is often explicitly closed for teaching establishments).
    The classification is confirmed by the national competent
    authority and notified to the institution.

    \item[Supervisory chain.] The national NIS2 supervisory
    authority oversees compliance. The national or governmental
    \gls{csirt} is the operational incident-response team
    for public-sector entities, and \loc{national-infosec-agency}
    publishes the State information-security policy that
    applies to public bodies.

    \item[Notification platform.] A single national
    incident-notification platform is typically the entry point for
    \gls{nis2}, \gls{gdpr} and CER notifications. The binding
    cadence is 24~hours (early warning) / 72~hours (incident
    notification) / one month (final report).

    \item[National strategy hook.] National cybersecurity strategies
    commonly identify cooperation between the State and the
    academic and research sphere as a priority objective. The
    contribution of this report to the regulator's protocols
    (Section~\ref{sec:advocacy}) is consistent with such a
    national objective.
\end{description}

\subsection{The architectural proposition in one page}

Five elements summarise the architecture proposed by this
report:

\begin{enumerate}[nosep]
    \item \textbf{Multi-zone segmentation (seven zones).}
    Administrative (high security), Teaching (standard),
    Research (flexible), \gls{sciencedmz} (high-throughput
    transfers), Personal \gls{byod}, IoT/Building (isolated),
    Sensitive enclave (compliance-critical). Each zone has its
    own security profile.
    \item \textbf{Trusted Activity Contexts.} Each activity
    operates in a coherent, validated environment with the data,
    tools, services and access rights appropriate to its
    purpose.
    \item \textbf{Federated governance.} A
    \emph{Digital Policy Commission} chaired by an elected
    academic, with representatives from departments and
    interdisciplinary centres. Responsibility charters formalise
    the autonomy of research units within the institutional
    minimum standards required by the applicable national NIS2
    transposition.
    \item \textbf{Researcher-centred AI cloud access.} The
    Principal Investigator classifies data and selects providers
    proportionately. Non-EU access operates on a cooperative
    declaration basis: the central IT service registers the
    access on PI request, with periodic \gls{dpo} audit
    (Section~\ref{sec:safeguards-ai}).
    \item \textbf{Continuous-improvement mechanism.} Per-zone
    compliance evidence aggregates into an institution-level
    maturity dashboard (an ISO/IEC~27005-aligned risk-analysis
    method or tool, such as a nationally provided open-source
    methodology), producing the annual maturity report
    submitted to \loc{nis2-supervisory-authority} where
    the transposition so requires.
\end{enumerate}

\smallskip
\noindent\textit{\footnotesize Technical implementation: each of
the five elements above is developed in depth in the standalone
Technical Reference~\cite{ng-tech-ref-2026} --- the Trusted
Activity Context in Chapter~2, the multi-zone segmentation in
Chapter~8, the federated governance in Chapter~3
(sovereignty grid), the AI access pattern in Chapter~7 (policy)
and Chapter~10 (data classification), the continuous-improvement
mechanism in Chapter~16 (conformance suite).}

\subsection{How to read the rest of this report}

The report is structured for several reading paths:

\begin{itemize}[nosep]
    \item \textbf{For decision-makers (\loc{academic-governing-bodies} and \loc{governing-board}):}
    Section~\ref{sec:stakeholders} presents per-audience
    messages; the accompanying executive summary (companion
    document) provides a condensed synthesis.
    \item \textbf{For legal review (CISO, DPO, legal service):}
    Section~\ref{sec:reglementation} contains the full regulatory
    analysis, with the methodology note at its entry;
    Section~\ref{sec:obligations-mapping} maps the
    NIS2 Article~21(2) obligations to the report's propositions;
    Section~\ref{sec:safeguards-ai} details the cloud AI access
    safeguards; Section~\ref{sec:advocacy} articulates the
    compliance posture in four positions.
    \item \textbf{For technical leads (IT services, faculty IT
    correspondents):} Section~\ref{sec:etat-art} surveys the
    state of the art; Section~\ref{sec:propositions-infra}
    details the proposed architecture and the seven zones;
    Section~\ref{sec:transition} sets out the transition plan.
\end{itemize}

% ============================================================================
\section{Introduction}
\label{sec:introduction}
% ============================================================================

\epigraph{``Those who would give up essential Liberty, to purchase a little temporary Safety, deserve neither Liberty nor Safety.''}{--- Benjamin Franklin, 1755~\cite{franklin1755}}

\subsection{Context and problem statement}

\textcolor{Crimson}{\textbf{A university requires administration, but is not a corporation;}}
university governance must optimise for freedom for ideas and persons,
physical or digital. The digital infrastructure should reflect this
fundamental distinction.

\paragraph{The liberty-first rule.}
\label{par:liberty-first}
This report applies one standing rule wherever freedom and constraint
collide: \keyword{constraint is never preferred to freedom by
default}, and the burden of proof always rests on the restriction.
The rule escalates through three levels. \emph{(i)~Design:} every
control is proportionate, documented, and limited to what is strictly
necessary; the default is freedom, the exception is justified.
\emph{(ii)~Contracts:} an engagement already entered into --- a vendor
agreement, an insurance policy --- that would limit an academic
freedom is, first of all, a contract to \keyword{renegotiate} so as to
maximise freedom in the academic context; freedom has a price, and the
governance bodies may knowingly choose to pay it (residual-risk
retention, an alternative counterparty, sector mutualisation) rather
than surrender the mission. \emph{(iii)~Regulation:} where the
constraint comes from law or regulation, the institution engages ---
collectively, with the national and international university
community --- in the influence-and-reform processes of
Section~\ref{sec:advocacy} (Position~C), so that legislators and
regulators do not impose regimes, insurance terms included, whose only
outcome would be the deprivation of academic digital freedoms.

\medskip
Universities are, by nature, open spaces dedicated to the creation
and transmission of knowledge. Their mission rests on three fundamental
pillars: \keyword{research}, which requires maximum freedom
of experimentation; \keyword{teaching}, which demands access to
diverse tools and the ability to configure varied pedagogical
environments; and \keyword{administration}, which requires the
protection of sensitive data (personal data, financial data,
human resources).

In a research-and-teaching university, these missions are typically
carried out by organisationally distinct entities: \emph{research-only units},
\emph{research-and-teaching departments}, and
\emph{teaching/training-only components}, each with different
infrastructure needs and different regulatory exposure
(see Section~\ref{sec:reglementation}).

Faced with the proliferation of cyber threats, some universities adopt
a \emph{centralised} security approach consisting of:
\begin{itemize}[nosep]
    \item Minimising the number of authorised software;
    \item Fully centralising systems administration;
    \item Restricting access rights uniformly for all
          staff;
    \item Prohibiting the installation of servers or local services in
          laboratories and departments.
\end{itemize}

While this approach may be justified for a purely
administrative environment, it \keyword{raises significant challenges} in a
university context for several reasons:

\begin{enumerate}[nosep]
    \item \textbf{It can hinder experimental research} in computer science,
    data science, artificial intelligence and any
    discipline requiring configurable infrastructures.
    \item \textbf{It can degrade the quality of teaching} across
    all faculties --- preventing technical courses from using
    realistic environments, and preventing non-technical
    disciplines from accessing AI tools, simulation platforms,
    and specialised digital resources.
    \item \textbf{It can create a false sense of security}: reducing
    the attack surface through total restriction risks concentrating
    vulnerabilities on a single point of failure.
    \item \textbf{It tends to drive \emph{shadow IT}}: researchers and
    instructors may circumvent restrictions by using unapproved
    solutions, potentially creating risks greater than those
    the policy claims to avoid~\cite{shadow-it-umbc}.
    \item \textbf{It hinders access to artificial intelligence},
    whose global cloud services (OpenAI, Google, AWS, HuggingFace)
    have become widely used across contemporary research and teaching.
\end{enumerate}

\subsection{Objectives of this report}

The objectives of this report are to:
\begin{enumerate}[nosep]
    \item Present a \textbf{state of the art} of university
    IT infrastructure architectures;
    \item Demonstrate through \textbf{concrete examples} that decentralised
    and segmented architectures are common across a range of
    research-intensive universities;
    \item Present the \textbf{normative frameworks} (NIST, ISO, COBIT)
    that support this approach;
    \item Propose \textbf{infrastructure and governance solutions}
    adapted to the specific needs of a university.
\end{enumerate}

\subsection{The Enterprise IT vs Scientific IT dichotomy}

The core of the problem lies in the uniform application of an
\emph{Enterprise IT} policy to a \emph{Research IT} environment.
As the \gls{sciencedmz} design pattern demonstrates, the performance and
security characteristics of enterprise traffic and scientific traffic
are fundamentally different~\cite{dart2013}. Applying administrative
restrictions to a compute cluster or a genomic sequencer
amounts to treating a scientific instrument like a secretarial workstation.

\subsection{Shadow IT: documented effects of restrictive policies}
\label{subsec:shadow-it}

A case study from the University of Maryland, Baltimore County (UMBC) revealed that
Shadow IT emerges to satisfy \emph{``unmet functional
needs''} and provide \emph{``innovation and flexibility''}~\cite{shadow-it-umbc}.
When central tools are unsuitable or too locked down,
researchers deploy their own servers under desks, use
personal cloud services for sensitive data, or create
unauthorised Wi-Fi access points. These systems escape all
oversight: they are neither patched, nor backed up, and are often
vulnerable~\cite{shadow-it-threatlocker}. \textbf{Overly restrictive central control can reduce visibility
into actual practices on the ground.}

This pattern is not unique to UMBC. A \gls{geant} survey of European
research and education institutions found that \emph{``in most
institutions, all official IT can have a shadow version, mainly
deployed by individual researchers or IT admins at the department
level''}~\cite{geant-shadow-it}. Gartner estimates that 41\% of
employees acquired or created technology outside IT's visibility in
2022, rising to a projected 75\% by
2027~\cite{gartner-shadow-it}. A systematic review of 107
publications on Shadow IT~\cite{klotz2019} identifies the root
causes as poor business-IT alignment, lack of agility, and
inadequate policies --- not malicious intent. The evidence
consistently shows that \textbf{prohibition does not eliminate
Shadow IT; it merely removes institutional visibility and
governance}, increasing rather than decreasing actual
risk~\cite{silic2014,mallmann2018}.

\subsection{Fundamental principle: dependability}
\label{subsec:dependability}

\keyword{Dependability} is defined by Avizienis et
al.~\cite{avizienis2004} as \emph{``the ability to deliver service
that can justifiably be trusted''}. Their foundational taxonomy
identifies six core \keyword{attributes} that any trustworthy system must
exhibit:

\begin{table}[H]
\centering
\caption{Dependability and security attributes (after~\cite{avizienis2004})}
\label{tab:dependability-attributes}
\small
\begin{tabularx}{0.92\textwidth}{l X}
\toprule
\textbf{Attribute} & \textbf{Definition} \\
\midrule
Availability    & Readiness for correct service \\
Reliability     & Continuity of correct service \\
Safety          & Absence of catastrophic consequences on the user and environment \\
Integrity       & Absence of improper system alterations \\
Maintainability & Ability to undergo modifications and repairs \\
Confidentiality & Absence of unauthorised disclosure of information \\
\bottomrule
\end{tabularx}
\end{table}

\noindent
The same taxonomy identifies four \keyword{means} for achieving
dependability: \emph{fault prevention}, \emph{fault tolerance},
\emph{fault removal}, and \emph{fault forecasting}. In a university
environment where experimentation, heterogeneity and openness are
structural, \textbf{fault tolerance} --- the ability to deliver correct
service in the presence of faults --- deserves particular emphasis
alongside fault prevention. Avizienis et al.\ present the four
means as complementary~\cite{avizienis2004, guelfi2007ftsystems}; the specific
characteristics of university environments (heterogeneity, \gls{byod},
academic freedom, experimentation) shift the balance toward
tolerance and detection rather than prevention alone.

\medskip
Two complementary mechanisms are central to fault tolerance in IT
architecture.

The first is \keyword{segmentation}: it answers the question
\emph{``Where can network flows circulate?''} by dividing the
infrastructure into distinct zones at the network level to
\textbf{limit the propagation} of incidents (lateral movement,
malware spread, cascading failures). Segmentation acts on the
\emph{blast radius} of an event. It is a well-established concept in
the standards literature, appearing as \emph{boundary protection}
(NIST SP~800-53 SC-7~\cite{nist-sp800-53}), \emph{network
segregation} (ISO~27002 A.8.22~\cite{iso27001}), and
\emph{segmentation} among the fourteen cyber-resiliency techniques
of NIST SP~800-160 Vol.~2~\cite{nist-sp800-160v2}. The second
mechanism --- which complements segmentation --- is developed in
the next subsection.

\subsection{Trusted Activity Contexts}
\label{sec:tac}

The second of the two mechanisms introduced above addresses a
complementary concern: not \emph{where} flows circulate, but
\emph{what validated environment} a user operates in for a given
activity. The standards literature
approaches this through several related but partial concepts:
\emph{isolation} (NIST SC-7(21)), \emph{privilege restriction} (NIST
SP~800-160 Vol.~2), \emph{attribute-based access control} (NIST
SP~800-162~\cite{nist-sp800-162}), and \emph{role-based session
activation} (\gls{rbac}~\cite{sandhu1996}). The practice literature
additionally uses terms such as \emph{compartmentalisation}
(Verizon DBIR~\cite{verizon-dbir2024}), \emph{enclaves} (NIST
SP~800-171~\cite{nist-171}), and \emph{security domains}
(CNSSI~4009). Each captures a facet of the problem, but none
defines a \textbf{complete, pre-assembled, activity-centric
environment} that brings together data, tools, services, and access
rights as a coherent unit for a specific mission.

To fill this gap, we introduce the concept of \keyword{Trusted
Activity Context} (\gls{tac}). A Trusted Activity Context is a coherent,
validated environment that assembles the data, tools, services, and
access rights required for a specific task, mission, or activity ---
such as a research project, a digital examination, an administrative
process, or a teaching session. Within a \gls{tac}, security policies,
access controls, logging, and isolation mechanisms are aligned with
the objectives and risk level of that specific activity.

\smallskip\noindent
\emph{Note:} the term \emph{Trusted Activity Context} should not be
confused with IBM DB2's identically named database connection
feature, nor with the Trusted Computing concepts of the Trusted
Computing Group (TCG). Here, ``trusted'' follows the dependability
definition of Avizienis et al.~\cite{avizienis2004}: an entity is
trusted if it delivers service that can justifiably be relied upon.

\medskip\noindent
The concept of \gls{tac} draws on and extends several established
frameworks: \gls{rbac} session-based role activation~\cite{sandhu1996},
Task-Based Authorization Controls (TBAC) which tie permissions to
task lifecycles~\cite{thomas-sandhu1997}, \gls{abac} which evaluates
contextual attributes for access decisions~\cite{nist-sp800-162},
and the contextual awareness technique of NIST SP~800-160
Vol.~2~\cite{nist-sp800-160v2}. However, none of these frameworks
defines a \emph{governance-level} abstraction that names and formalises
the combination of data, tools, services, and access rights as a
coherent, auditable unit tied to a specific university activity.
Cloud workspace products (AWS WorkSpaces, Azure Virtual Desktop) and
container orchestration (Kubernetes namespaces with \gls{rbac}) implement
aspects of this pattern operationally; \gls{tac}s provide the
\textbf{conceptual unification} needed for university governance,
elevating from per-request access decisions to
\textbf{per-activity environment composition}.

\medskip\noindent
\gls{tac}s rest on four fundamental principles:

\begin{description}[style=nextline, leftmargin=1.5cm, nosep]
    \item[\textsc{Trustworthiness}] Each context is technically
    controlled, monitored, and bounded so that users can be granted a
    high degree of autonomy within it.

    \item[\textsc{Dependability}] Failures, misconfigurations, or
    security incidents occurring within one context remain contained
    and cannot compromise the core infrastructure or other contexts.

    \item[\textsc{User-centricity}] The digital environment adapts to
    the user's current role and activity rather than imposing a rigid
    and uniform security model on all users and all situations.

    \item[\textsc{Observability}] Every action within a \gls{tac} is logged
    and auditable, ensuring that autonomy is accompanied by
    accountability.
\end{description}

\noindent
\gls{tac}s transform isolation from a restrictive mechanism into an
\textbf{enabling} one. Isolation is no longer primarily a way to
restrict users, but a way to safely grant autonomy. By ensuring that
each activity takes place within a controlled and isolated
environment, the system protects its core functions while allowing
experimentation, research, teaching, and collaboration to take place
with a high degree of freedom.

\medskip\noindent
A mature university architecture thus combines a \emph{relatively
stable segmentation} (network zones) with multiple \emph{dynamic
Trusted Activity Contexts} that operate within those zones,
adapting to the roles, missions, and risk levels of each activity.

\medskip
The forestry analogy is illuminating: a forest
is not made safer by planting trees in a single continuous mass --- on the
contrary, it is the \emph{creation} of firebreaks that
prevents a blaze from spreading to the entire woodland.

\begin{tcolorbox}[
    colback=NavyBlue!6!white,
    colframe=NavyBlue!80!black,
    boxrule=0.8pt,
    arc=2pt,
    left=10pt, right=10pt, top=8pt, bottom=8pt,
    title={\textbf{Guiding principle for a university information system}}
]
\itshape
A university information system must not attempt to prevent all errors,
experiments, or unforeseen uses; it must prevent such events from
compromising the institution as a whole.

\medskip
\upshape
\textbf{In short:} local freedom, incident confinement, protection of
the core, global resilience.
\end{tcolorbox}

\smallskip\noindent
\emph{Note for technical readers.} A non-binding sketch of how the
\gls{tac} concept might be assembled from current open-source and
commercial building blocks --- end-to-end sequence, endpoint
operating-system support matrix, reference architecture by layer, and
a wave-based migration playbook --- is developed in the standalone
companion document~\cite{ng-tech-ref-2026}, with the explicit caveat
that the proposal would benefit from verification against the IT
department's actual environment.

\medskip\noindent
These are precisely the principles that major universities apply,
as the following sections demonstrate.

% ============================================================================
\section{European regulatory context}
\label{sec:reglementation}
% ============================================================================

\epigraph{``Useless laws weaken the necessary laws.''}{--- Montesquieu, \textit{The Spirit of the Laws}~\cite{montesquieu1748}}

\subsection*{Methodology note}

The regulatory analysis presented in this section is grounded in
the publicly available legislative texts (the NIS2
Directive~\cite{nis2} and its national transposition, the
Critical Entities Resilience (CER) Directive~(EU)~2022/2557 and
its national transposition, the
\gls{gdpr}~\cite{gdpr}, the European AI Act~\cite{eu-ai-act}, and
for digital-sector entities, Commission Implementing
Regulation~(EU)~2024/2690~\cite{eu-regulation-2024-2690}),
the binding case law (CJEU \emph{Schrems~II}~\cite{schrems2020}),
and the national instruments (the national public-sector
information-security policy and \loc{national-cyber-strategy}\loc{cite-cyber-strategy}'s
academic-cooperation objective) cited throughout. The
analysis is produced by the contributors of this report acting in
their academic-community capacity; it has not been formally
validated by the institution's legal service. The contributors
welcome formal legal review and remain available to discuss any
specific correction that such review may identify. The reference
shelf and citation chain enable any reader to verify each claim
against its primary source.

\subsection{GDPR Article 32: a \emph{risk-based} obligation}

Article~32 of the \gls{gdpr}~\cite{gdpr} requires data controllers
to implement technical and organisational measures
\emph{``appropriate''} to ensure a level of security \emph{``adapted
to the risk''}, taking into account the \emph{``state of the art''}, costs,
context and the likelihood/severity of risks. It explicitly
mentions:

\begin{itemize}[nosep]
    \item The ability to ensure the \textbf{confidentiality, integrity,
    availability} and \keyword{resilience} of systems;
    \item The ability to \textbf{restore availability} within
    appropriate timeframes;
    \item A process for \textbf{regularly testing and evaluating}
    the effectiveness of measures.
\end{itemize}

The \gls{edpb} (\emph{European Data Protection Board})~\cite{edpb-guidelines}
and \gls{enisa}~\cite{enisa-risk} confirm that Article~32 is based on the
\emph{``appropriate to the risk''} standard and that measures must be
adapted to the context.

\textbf{Implication for universities:} it is legally and
operationally defensible to operate \textbf{multiple security
domains} (administrative, research, teaching) with differentiated
freedoms, provided that each domain has:
\begin{enumerate}[nosep]
    \item a documented risk assessment;
    \item defined and traceable responsibilities;
    \item controls proportionate to the sensitivity of the data;
    \item evidence of regular effectiveness testing.
\end{enumerate}

\subsection{NIS2 Directive: the institution as an essential entity}
\label{sec:nis2}

Where the institution's jurisdiction has transposed the \gls{nis2}
Directive~\cite{nis2}, a public research-and-teaching institution \loc{essential-entity-modality}
as an \keyword{essential entity}\loc{nis2-status-conditional-tail}\loc{essential-entity-classification-basis}. This classification is
independent of the institution's research activities and applies to the
institution as a whole.

Such a classification rests on three cumulative legal pegs that the
national transposition typically reproduces:
the definition of the \keyword{public-administration entity}
through four cumulative criteria (public-interest mission, legal
personality, public funding or State control, decision-making power
affecting cross-border movements); the clause under which a
public-administration entity is essential \emph{regardless of size};
and the provision which excludes teaching establishments
from the ``research organisation'' track that would otherwise apply
the lighter important-entity regime.

This report recommends that the institution accept this classification,
where it applies, as the binding operational basis of its \gls{nis2}
compliance. Acceptance of the outcome does not entail suspension of
intellectual scrutiny over the classification protocol itself: as a
public knowledge institution hosting \loc{cs-research-unit} whose
research portfolio covers information-systems security, the institution
retains and is invited by this report to exercise a duty to contribute
technical and methodological input to the supervisory authorities'
classification practice (see Section~\ref{sec:advocacy}).

This classification has three direct consequences:

\begin{enumerate}
    \item \textbf{Holistic perimeter.} \gls{nis2} obligations apply to
    the institution in its entirety, with no possibility to narrow the
    scope to specific services, faculties or units --- a structural
    difference from NIS1.
    \item \textbf{Higher level of obligations.} Essential entities are
    subject to proactive (\emph{ex-ante}) supervision rather than the
    reactive (\emph{ex-post}) oversight applicable to important
    entities, with heavier sanctions (Article~34); the Article~23
    reporting cadence itself is the same for both classes --- the
    differential lies in supervision and sanctions.
    \item \textbf{Continuous compliance evidence.} Where the national
    transposition so provides (as in the worked example), the
    institution is required to submit a regular compliance and
    maturity report to the competent authorities --- a national
    mechanism reflecting the directive's \emph{ex-ante} supervision
    powers (Article~32) --- which presupposes governance, audit and
    incident-management evidence at the level of the entire
    institution.
\end{enumerate}

\textbf{Article~21 proportionality.} A holistic perimeter does not
require uniform controls. \gls{nis2} Article~21(1) requires
\emph{appropriate and proportionate} technical, operational and
organisational measures, taking into account the state of the art,
costs of implementation, and the entity's exposure to risk. The
multi-zone architecture proposed in
Section~\ref{sec:propositions-infra} is precisely the mechanism through
which Article~21 proportionality is implemented across a heterogeneous
public-research perimeter: it operationalises \gls{nis2} compliance, it
does not narrow the scope.

\subsection{National institutional landscape}
\label{sec:lu-landscape}

The European-level framework set out above is operationalised at
national level through a jurisdiction-specific institutional and
procedural ecosystem. The present subsection anchors the five elements
directly relevant to the institution --- stated in role-based terms ---
so that the proposed architecture is opposable point by point to the
national bodies that supervise it. (The worked instantiation for a
given regime is given in the Regulatory Applicability annex.)

\subsubsection{Supervisory authority and incident-response chain}

Three categories of national body typically hold direct,
non-overlapping responsibilities towards the institution:

\begin{description}[nosep]
    \item[The NIS2 supervisory authority.] The competent authority
    designated under the national transposition of
    \gls{nis2}; the body that issues the institution's
    essential-entity classification and to which the annual compliance
    and maturity report is addressed.
    \item[The national / governmental \gls{csirt}.] The
    incident-response team for public-sector entities, namely
    \loc{national-gov-csirt}\loc{cite-gov-csirt}. As a public
    institution, the institution reports incidents primarily through
    this team.
    \item[The national single point of contact.] The body responsible
    for cross-border NIS2 cooperation and major-incident escalation;
    it commonly hosts both \loc{national-gov-csirt} and \loc{national-infosec-agency}.
\end{description}

\noindent Two further categories of body may exist for ecosystem
disambiguation: \loc{national-private-sector-csirt}\loc{cite-circl} (relevant when third-party private
partners are involved), and \loc{financial-sector-supervisor}\loc{cite-cssf} (competent
only for the financial sector and not for the institution outside
specific financial projects).

\subsubsection{Notification platform}

Many jurisdictions operate a single national incident-notification and
risk-management platform that consolidates notifications across NIS,
\gls{nis2}, EECC, \gls{gdpr} and CER directives, operated by the
supervisory authority in partnership with the national cybersecurity
bodies and \loc{national-dpa}. The institution,
as an essential entity, discharges its NIS2 incident-notification
obligations through such a platform on the binding 24-hour / 72-hour /
one-month cadence (\gls{nis2} Article~23) and submits its annual maturity
report through the same channel.

\subsubsection{National strategy hook}

National cybersecurity strategies commonly identify the strengthening
of cooperation with the \emph{academic and research sphere} as a
national priority. Such an objective is the direct national-level
political hook for the present report's positioning: the institution's
contribution to the regulator's protocols (Section~\ref{sec:advocacy})
is consistent with it.

\subsubsection{Concomitant CER framework}

\Loc{cer-national-transposition}\loc{cite-cer-transposition}, transposing the Critical Entities
Resilience (CER) Directive~(EU)~2022/2557, typically provides that every entity
designated as critical under the CER framework is automatically
essential under NIS2. \loc{cer-designation-status}; the cross-reference is provided for
regulatory-neighbourhood completeness.

Commission Implementing
Regulation~(EU)~2024/2690~\cite{eu-regulation-2024-2690} of
17~October~2024 details the ten obligation families of NIS2
Article~21(2) for the digital-sector entity classes within its
Article~1 scope; for other essential or important entities ---
universities included --- it is not directly binding but serves as
persuasive state-of-the-art guidance, the binding baseline being set
by the national transposition.

\subsubsection{Public-sector reference framework}

The public-sector information-security baseline is typically governed
by \loc{national-public-sector-isp}\loc{cite-public-sector-isp}, developed
by \loc{national-infosec-agency}\loc{cite-infosec-agency} and applicable to State
administrations. Such a baseline establishes the objectives of
confidentiality, integrity, availability and risk-based
security-measure adoption. The multi-zone architecture
proposed in Section~\ref{sec:propositions-infra} is positioned as a
compatible operational implementation that adapts that baseline
to the heterogeneous public-research perimeter of the institution.

\subsection{Operational obligations and proposition-mapping}
\label{sec:obligations-mapping}

The defensibility of the present report rests on the demonstration
that every obligation imposed by the national NIS2
transposition on essential entities --- mirroring NIS2 Article~21(2)
--- is already addressed by the architectural propositions developed in
Section~\ref{sec:propositions}. Two mapping tables make this visible.

\begin{table}[H]
\centering
\caption{NIS2 Article~21(2) ten-family obligations mapped to the report's propositions}
\label{tab:art12-mapping}
\small
\begin{tabularx}{\textwidth}{%
  >{\hsize=1.05\hsize\linewidth=\hsize}X%
  >{\hsize=0.75\hsize\linewidth=\hsize}X%
  >{\hsize=1.20\hsize\linewidth=\hsize}X%
}
\toprule
\textbf{Art.~21(2) family} & \textbf{Reference} &
    \textbf{Propositions covering} \\
\midrule
1. Risk-analysis and SI security policy &
    Art.~21(2)(a) & P1 (risk-based approach),
    P14 (per-zone evidence) \\
\addlinespace
2. Incident handling & Art.~21(2)(b) &
    P12 (security communities), P14 (evidence) \\
\addlinespace
3. Business continuity (backup, recovery, crisis) &
    Art.~21(2)(c) & P9 (\gls{sciencedmz}),
    P11 (secure alternatives) \\
\addlinespace
4. Supply-chain security & Art.~21(2)(d) &
    P6 (open-source by default), P7 (enclave charters) \\
\addlinespace
5. SI acquisition, development, maintenance security &
    Art.~21(2)(e) & P11 (secure alternatives),
    P6 (open-source) \\
\addlinespace
6. Effectiveness measurement &
    Art.~21(2)(f) & P14 (per-zone evidence +
    maturity dashboard) \\
\addlinespace
7. Cyber-hygiene practices and training &
    Art.~21(2)(g) & P5 (Digital Policy Commission),
    training in P10 / P12 \\
\addlinespace
8. Cryptography policies & Art.~21(2)(h) &
    P11 (secure alternatives), implicit in P3 (BYOD \gls{mfa}) \\
\addlinespace
9. HR security, access control, asset management &
    Art.~21(2)(i) & P3 (BYOD), P7 (charters),
    P8 (autonomy with minimums), P10 (federated identity) \\
\addlinespace
10. \gls{mfa} and secured communications &
    Art.~21(2)(j) & P13 (Zero Trust + \gls{mfa}),
    P10 (federated identity) \\
\bottomrule
\end{tabularx}
\end{table}

\noindent For the digital-sector entity classes within its
Article~1 scope, the technical baselines for each family are
detailed in Commission Implementing Regulation
(EU)~2024/2690~\cite{eu-regulation-2024-2690}; for other entities
they serve as persuasive state-of-the-art guidance.

\begin{table}[H]
\centering
\caption{Essential-entity-specific operational obligations}
\label{tab:obligations-essentielles}
\small
\begin{tabularx}{\textwidth}{X l X X}
\toprule
\textbf{Obligation} & \textbf{Reference} &
    \textbf{Current status} & \textbf{Planned action} \\
\midrule
Auto-registration with the supervisory authority &
    National NIS2 transposition &
    Pending & Initial declaration (transition plan, step~0) \\
\addlinespace
Notification of measures to the supervisory authority &
    National NIS2 transposition &
    Architecture produces the evidence (P14) & Annual maturity
    report via the national notification platform \\
\addlinespace
Ex-ante supervision & \gls{nis2} essential-entity regime &
    Architecture is audit-ready & Internal-audit cycle feeding
    the maturity report \\
\addlinespace
Incident notification 24\,h / 72\,h / 1~month &
    \gls{nis2} Art.~23 &
    Section~\ref{sec:lu-landscape} introduces the platform &
    Incident-response procedure integrated with the platform's cadence \\
\addlinespace
Personal liability of management bodies &
    \gls{nis2} Art.~20 & Implicit in governance & Documented in
    the executive-training programme
    (Section~\ref{sec:exec-training}) \\
\addlinespace
Executive training of management bodies &
    \gls{nis2} Art.~20(2) & Implicit in Digital Policy Commission
    role & Programme via Digital Policy Commission
    (Section~\ref{sec:exec-training}) \\
\bottomrule
\end{tabularx}
\end{table}

\subsection{European AI Act and access to AI cloud services}

The European regulation on artificial intelligence (AI
Act, 2024)~\cite{eu-ai-act} introduces differentiated obligations
according to the risk levels of AI systems. For universities, the issue
is twofold:

\begin{enumerate}
    \item \textbf{AI research:} researchers must be able to
    freely access APIs and cloud services from global providers
    (OpenAI, Anthropic, Google, Meta, HuggingFace, AWS, Azure) to
    train, evaluate and compare models. Some of these services
    are not hosted in the EU and are not fully \gls{gdpr}
    compliant;
    \item \textbf{AI teaching:} advanced courses require
    access to these same platforms for hands-on exercises on
    foundation models.
\end{enumerate}

The AI Act~\cite{eu-ai-act} provides two distinct provisions relevant
to research:
\begin{itemize}[nosep]
    \item \textbf{Article~2(6):} exempts AI systems
    \emph{specifically developed and put into service for the sole purpose}
    of scientific R\&D. This covers bespoke research tools built by
    researchers, \textbf{not} the use of commercial AI services
    (ChatGPT, Claude, Gemini) as research instruments;
    \item \textbf{Article~2(8):} exempts research, testing and development
    activities on AI systems \emph{prior to their being placed on the
    market}.
\end{itemize}
For researchers \emph{using} commercial AI APIs, the relevant
framework is the AI Act's \emph{risk-based classification}: the
applicable obligations depend on the concrete use case and the user's
role under the Act; many research uses fall outside the high-risk
categories. This risk classification ---
not a blanket research exemption --- is the defensible legal basis for
differentiated AI cloud access in the research zone (see
Section~\ref{sec:ia}).

\subsection{Regulatory summary}

\begin{table}[H]
\centering
\caption{European regulatory framework and architectural implications}
\label{tab:reglementation}
\small
\begin{tabularx}{\textwidth}{l X X}
\toprule
\textbf{Regulation} & \textbf{Key requirement} &
    \textbf{Architectural implication} \\
\midrule
\gls{gdpr} Art. 32 & Measures appropriate to the risk, resilience, testing &
    Multi-zone architecture with per-zone assessment \\
\addlinespace
NIS2 (essential entity) & Holistic perimeter, proactive supervision,
    mandatory audits, annual maturity report; Art.~21 proportionate measures &
    Article~21 proportionality via multi-zone segmentation; enclave
    registry, incident playbooks, formalised federated governance \\
\addlinespace
NIS2 (operational cadence) & 24\,h/72\,h/1-month incident notification
    via the national notification platform + annual maturity report to
    the supervisory authority & Incident-response procedure integrated
    with the national platform; institution-level maturity dashboard
    producing the annual report \\
\addlinespace
AI Act & Risk-based classification; obligations by risk level &
    Research zone with managed, differentiated AI cloud access \\
\addlinespace
\gls{gdpr} (transfers) & Standard contractual clauses for transfers outside the EU &
    Legal framework for access to US clouds (research) \\
\bottomrule
\end{tabularx}
\end{table}

\subsection{Compliance posture and knowledge-institution responsibility}
\label{sec:advocacy}

The propositions developed in this report cluster into four
labelled positions, each named to clarify the nature of the
contribution under the corresponding regulatory framework. The
labelling is intended to make the dossier opposable to the
supervisory authorities point by point and to distinguish
operational compliance from areas where the report proposes a
substantive interpretation or contributes to the regulator's
protocol-improvement processes.

\paragraph{Note on tone.}
Throughout the four positions, references to the classification
process, to the regulator's protocols, and to areas of ongoing
methodological debate use the language of constructive
contribution, not of dispute. The relationship the report proposes
between the institution and the supervisory authorities is
collaborative: the binding character of the regulatory framework
is acknowledged in every relevant section; the institutional duty
of intellectual contribution is acknowledged in parallel, without
prejudice to the former.

\paragraph{Position A --- Compliance.}
The majority of the report's propositions (P1, P3, P5, P6, P7,
P9, P10, P11, P12, P13, P15) are unambiguously compliant under
the national NIS2 transposition, given the
architecture and operational mapping developed in this report.
These propositions require no further intellectual treatment
beyond the operational implementation roadmap (see
Section~\ref{sec:transition}); they form the compliance baseline
of the dossier.

\paragraph{Position B --- Documented interpretation.}
Three areas in which the report's interpretation is documented
with its supporting reasoning so that it can be examined on its
merits.

\subparagraph{B.1 --- Cross-border decision-making power (public-administration criterion).}
This report supports that the institution satisfies the fourth
criterion of the public-administration entity definition through
diploma issuance, ECTS recognition, and qualification certification.
The interpretation is reasonable but novel; this report recommends
that interpretive guidance be requested from the national supervisory
authority and \loc{national-spoc-authority}\loc{cite-spoc} as part of the
auto-registration dialogue
(Section~\ref{sec:obligations-mapping}).

\subparagraph{B.2 --- Cloud AI access for research under \emph{Schrems~II} + AI Act.}
The cumulative regime articulated in
Section~\ref{sec:safeguards-ai} is the strictest contemporary
interpretation of CJEU \emph{Schrems~II}~\cite{schrems2020} in
combination with the AI Act~\cite{eu-ai-act}. This report
indicates that the TIA template and the data-classification
table could be made available, as contributing artefacts, to the
\gls{edpb} / \gls{edps} doctrinal effort on research-context transfers, to
the extent that institutional decisions and resources allow.

\subparagraph{B.3 --- NIS2 Art.~21 proportionality applied to a multi-mission perimeter.}
The seven-zone segmentation as NIS2 Article~21 implementation
(Section~\ref{sec:nis2}) is the operational thesis of this
dossier. Prior practice in the public sector has frequently not
formalised such a multi-zone approach; this report proposes it as
a model and indicates that the supporting evidence (annual
maturity report structure, risk-analysis-method instances) could
be shared with peer public-administration entities, subject to
institutional decision.

\paragraph{Position C --- Influence and reform.}
Four areas in which this report identifies a legitimate channel
for the institution, as a knowledge institution, to inform the
legislator and the regulator about the operational consequences
of current rules.

\subparagraph{C.1 --- Exclusion of teaching establishments from the ``research organisation'' track.}
Where the national transposition closes the
research-organisation route for universities, it places them under
the heaviest regime (essential entity via public-administration)
by default. This report indicates that the institution is well
placed to contribute to a post-transposition review of whether
this allocation is operationally optimal --- without advocating
for a downgrading of the present classification.

\subparagraph{C.2 --- ``Essential regardless of size'' rule.}
Where this rule applies it is binding, and this report supports its
full operational acceptance. The report indicates that the institution
could contribute, as a knowledge institution, to evidence-based
discussion at EU level on whether tiered obligations within the
essential category would improve overall NIS2 effectiveness ---
without advocating for any individual derogation.

\subparagraph{C.3 --- Right of the regulated entity to request review of its own classification.}
The national transposition, read in
conjunction with the applicable national administrative-review
framework, does not preclude the regulated
entity from submitting evidence-based observations on the
classification methodology applied to its case. This report
supports that the institution does not currently exercise this
right with respect to the substance of the decision (the
essential-entity status is accepted as binding). The report
indicates that this possibility could be reserved should
methodological refinements emerge --- particularly around the
interpretation of the cross-border
decision-making criterion --- that materially affect the
classification of comparable institutions.

\subparagraph{C.4 --- AI Act research-support principle and simplified notification.}
AI Act Recital~25 and the innovation-support measures of Chapter~VI
(Articles~57--63)~\cite{eu-ai-act} call for
research and innovation activities to be supported and not unduly
restricted. This report supports the proposition that the relevant
legislator formalise a simplified notification framework for transfers of
pseudonymised personal data in research contexts, in addition to
the AI Act Articles~2(6) and~2(8) exemptions for bespoke research
tools and pre-market R\&D. This report indicates that the
institution could contribute to a national or EU consultation on
this framework, consistent with \loc{national-cyber-strategy}'s
academic-cooperation objective and with the cooperative declaration
model articulated in Section~\ref{sec:safeguards-ai}.

\paragraph{Position D --- Methodological contribution to the supervisory authority.}
As a public knowledge institution hosting \loc{cs-research-unit}
whose research portfolio covers information-systems
security, the institution holds an intellectual and moral
responsibility that exceeds mere compliance: a responsibility to
share rigorous technical and scientific input with the bodies
that design and operate the national cybersecurity-regulation
protocols (the national supervisory authority, the national single
point of contact and governmental CSIRT, and \loc{national-infosec-agency}).
The present report is itself the first concrete expression of
this contribution; the domains listed below are offered without
commitment of resources or specific deliverables, on a
case-by-case basis subject to the formal request of the relevant
authority.

\subparagraph{D.1 --- Stress-testing of classification criteria.}
This report indicates that the institution could provide, on
request or unsolicited, a peer-reviewed analysis of the
four-criteria interpretation of the public-administration-entity
definition, in particular the cross-border decision-making
criterion, with comparative analysis across EU public-research
entities.

\subparagraph{D.2 --- Open-source methodology and tooling.}
This report indicates that the institution could share, on
request, the TIA template (Section~\ref{sec:safeguards-ai}), the
data-classification table and the maturity-report
structure adopted by the present report, as reference artefacts
that the regulator and peer entities may adapt --- subject to
institutional decisions and resource availability.

\subparagraph{D.3 --- Consultative engagement.}
This report indicates that the computer-science department's
expertise could be made available to the national supervisory
authority's, the single point of contact's and the
information-security agency's consultative processes, including
standardisation working groups, post-transposition reviews of
the national NIS2 transposition, and EDPB / ENISA working drafts
when academic input is solicited.

\subparagraph{D.4 --- Academic-regulator dialogue artefacts.}
The annual maturity report submitted to the
national supervisory authority pursuant to the national
transposition is, on the institution's side, a compliance
artefact. This report indicates that the institution could
voluntarily complement it with a short methodological annex
sharing observations and proposed protocol refinements arising
from the year's operational experience --- subject to
institutional decisions and resource availability.

\subparagraph{D.5 --- Research-context interpretation of the GDPR accountability principle (\emph{Schrems~II} application).}
This report proposes a cooperative researcher-centred model for
personal-data transfers in research contexts: PI data
classification, institutional register maintained by the central
IT service on PI request, periodic \gls{dpo} audit, with no
per-access gatekeeping. The model operationalises \gls{gdpr}
Articles~24(1) and~5(2) in a way that preserves international
scientific collaboration as an academic right. This report
indicates that the institution could share its register template
and audit methodology with the national supervisory authority,
\loc{national-dpa} and
EDPB working groups, as a contribution to research-context
regulatory practice --- without commitment of resources beyond
the present proposal.

% ============================================================================
\section{Specificities of universities regarding IS}
\label{sec:specificites}
% ============================================================================

\epigraph{``The university is not one community but several~[\ldots]. Its edges are fuzzy~[\ldots]. A community should have a soul, a single animating principle; the multiversity has several.''}{--- Clark Kerr, \textit{The Uses of the University}~\cite{kerr1963}}

\subsection{The university as an open space}

Universities fundamentally differ from businesses in:

\begin{description}
    \item[User diversity:] students (transient,
    personal equipment), academic staff (demanding autonomy),
    administrative staff, visiting researchers, international external
    collaborators, alumni.
    \item[Equipment diversity:] managed workstations, \gls{byod}
    (\emph{Bring Your Own Device}), research instruments (often running
    legacy operating systems), IoT sensors, \gls{hpc} clusters, laboratory
    servers.
    \item[Culture of openness:] academic freedom implies that
    researchers can choose their tools, test configurations,
    deploy prototypes and collaborate internationally without
    bureaucratic obstacles.
    \item[Scientific openness:] the \emph{Open Science}~\cite{unesco-open-science,
    open_knowledge_foundation_open_2015, berners-lee_linked_2006} and
    \emph{Open Data}~\cite{fair2016} movements push towards fewer
    restrictions on sharing and access to data and tools.
    \item[AI as an everyday tool:] in 2025--2026, generative
    artificial intelligence has become an essential working tool
    for research and teaching, requiring access to
    global cloud services.
\end{description}

\subsection{Differentiated needs by mission}

The table~\ref{tab:besoins} summarises the contrasting needs of the three
university missions.

\begin{table}[H]
\centering
\caption{Differentiated needs by university mission}
\label{tab:besoins}
\small
\begin{tabularx}{\textwidth}{l X X X}
\toprule
\textbf{Criterion} & \textbf{Administration} & \textbf{Teaching} & \textbf{Research} \\
\midrule
Traffic profiles & Predictable, web/email & Moderate, LMS/labs & Sustained multi-Gbps flows \\
\addlinespace
Security model & Stateful firewall, proxy & Standard controls & \gls{acl}, \gls{ids}, monitoring \\
\addlinespace
Latency tolerance & Moderate & Moderate & Very low (HPC) \\
\addlinespace
Bandwidth & 1--10 Gbps & 1--10 Gbps & 10--400 Gbps \\
\addlinespace
Ext. connectivity & Filtered Internet & Filtered Internet & Direct to NRENs \\
\addlinespace
AI cloud access & Not required & AI APIs for courses & \textbf{Unrestricted, including outside EU} \\
\addlinespace
Data sensitivity & Strict \gls{gdpr}, financial & Moderate & Variable (open to classified) \\
\addlinespace
Population & Admin. staff & Students, instructors & Researchers, collaborators \\
\addlinespace
Required autonomy & Low & Medium & \textbf{High (root access)} \\
\addlinespace
BYOD & Limited & Dominant (students) & Frequent (researchers) \\
\bottomrule
\end{tabularx}
\end{table}

\noindent\emph{Note:} in a research-and-teaching university,
\emph{``Research''} encompasses both research-only units and the research
activities of research-and-teaching departments; \emph{``Teaching''}
encompasses both teaching/training-only components and the teaching
activities of research-and-teaching departments.
Research-and-teaching departments thus span two zones.

\subsection{Effects of uniform restriction policies}

The approach of reducing the attack surface through uniform maximum restriction
suffers from several contradictions:

\begin{enumerate}
    \item \textbf{Single point of failure:} a centralised,
    undifferentiated infrastructure creates a \emph{single point of
    failure} --- if the central system is compromised,
    \emph{everything} is compromised.
    \item \textbf{Shadow IT:} excessive restrictions push users toward
    unmanaged solutions (personal clouds, servers under desks, personal
    \gls{vpn}s), creating major security
    blind spots~\cite{shadow-it-umbc,shadow-it-threatlocker}.
    \item \textbf{Loss of skills:} computer science and cybersecurity
    researchers can no longer train students in system realities if they
    do not have access to administrable environments.
    \item \textbf{Brain drain:} top-level researchers leave institutions
    where restrictions prevent them from working effectively, in favour of
    more flexible universities.
\end{enumerate}

\subsection{Staff BYOD: a structural need}

Staff \gls{byod} (\emph{Bring Your Own Device}) --- and not only
student \gls{byod} --- is an unavoidable reality:

\begin{itemize}[nosep]
    \item Faculty members often use a personal laptop as their primary
    work tool;
    \item Internationally mobile researchers must be able to connect
    from any device;
    \item Prohibiting \gls{byod} amounts to pushing toward Shadow IT
    (connections via personal hotspots, work on unmanaged clouds).
\end{itemize}

The solution is a \textbf{dedicated \gls{byod} subnet} offering:
\begin{enumerate}[nosep]
    \item Full Internet access;
    \item \keyword{Progressive interoperability} with other networks ---
    that is, a graduated model in which the scope of accessible
    institutional resources increases with the level of authentication
    and device posture verification:
    \begin{itemize}[nosep]
        \item Without authentication: Internet only (guest zone);
        \item With \gls{mfa} authentication: access to teaching/research zone
        resources matching the user's profile;
        \item With compliant device (verified posture): extended access
        including administrative resources via secure \gls{vdi}.
    \end{itemize}
\end{enumerate}

% ============================================================================
\section{State of the art: university infrastructure architectures}
\label{sec:etat-art}
% ============================================================================

\epigraph{``The good fighters of old first put themselves beyond the possibility of defeat, and then waited for an opportunity of defeating the enemy.''}{--- Sun Tzu, \textit{The Art of War}~\cite{suntzu-giles1910}}

\subsection{Protected high-throughput research-data paths: the Science DMZ pattern}
\label{subsec:science-dmz}

Research traffic that the general stateful-filtering perimeter penalises
needs a dedicated, protected high-throughput path. Several realisations
meet this function --- \gls{nren}-operated dedicated circuits and
lightpaths (\gls{geant}), experiment-operated overlays (CERN
LHCOPN/LHCONE), and tuned Data-Transfer-Node fabrics --- of which the
best-documented and most-instrumented is the \emph{Science DMZ} pattern,
a network architecture developed by ESnet (\emph{Energy Sciences
Network}) of the U.S. Department of Energy specifically for the needs of
high-performance scientific research~\cite{dart2013,geant,cern-lhcone}.
The remainder of this section describes the Science DMZ as the canonical
instantiation; the same function is realised elsewhere by the NREN- and
experiment-operated paths just named.

\subsubsection{Principles}

The \gls{sciencedmz} is a portion of the network, built at the campus perimeter,
whose equipment, configuration, and security policies are
\emph{optimised for scientific applications} rather than for enterprise
systems.

\begin{figure}[H]
\centering
\begin{tikzpicture}[
    node distance=1.2cm and 2cm,
    box/.style={rectangle, draw, rounded corners, minimum width=2.8cm,
                minimum height=0.8cm, font=\small, align=center},
    bigbox/.style={rectangle, draw, rounded corners, minimum width=3.2cm,
                   minimum height=0.8cm, font=\small, align=center},
    cloud/.style={ellipse, draw, minimum width=3.5cm, minimum height=1cm,
                  font=\small},
    arr/.style={-{Stealth[length=3mm]}, thick}
]
    % Internet
    \node[cloud, fill=blue!10] (inet) {Internet / NREN};

    % Border router
    \node[box, below=of inet, fill=gray!20] (border) {Border router};

    % Two branches
    \node[bigbox, below left=1.5cm and 1.5cm of border, fill=red!10]
        (fw) {Enterprise\\firewall};
    \node[bigbox, below right=1.5cm and 1.5cm of border, fill=green!10]
        (sdmz) {Science DMZ};

    % Enterprise
    \node[box, below=of fw, fill=red!5] (admin) {Administrative\\network};
    \node[box, below=of admin, fill=red!5] (email) {Email, Web,\\SIS, ERP};

    % Science DMZ components
    \node[box, below=of sdmz, fill=green!5] (dtn) {DTN\\(\emph{Data Transfer Nodes})};
    \node[box, left=0.5cm of dtn, fill=yellow!10] (perf)
        {perfSONAR};
    \node[box, below=of dtn, fill=green!5] (hpc) {HPC / Research\\storage};
    \node[box, right=0.5cm of dtn, fill=orange!10] (ids) {Zeek IDS};

    % Arrows
    \draw[arr] (inet) -- (border);
    \draw[arr] (border) -- (fw);
    \draw[arr] (border) -- (sdmz);
    \draw[arr] (fw) -- (admin);
    \draw[arr] (admin) -- (email);
    \draw[arr] (sdmz) -- (dtn);
    \draw[arr] (dtn) -- (hpc);
    \draw[arr] (sdmz) -- (perf);
    \draw[arr] (sdmz) -- (ids);
\end{tikzpicture}
\caption{Science DMZ architecture: separation of research and
administrative traffic}
\label{fig:science-dmz}
\end{figure}

\subsubsection{Key components}

\begin{description}
    \item[Data Transfer Nodes (DTN):] servers dedicated to massive data
    transfers, using Globus~\cite{globus} (HTTPS-based managed transfer),
    optimised for high-throughput flows.
    \item[Security through ACLs:] while modern next-generation firewalls
    advertise threat-prevention throughput in the 100-Gbps class,
    inspecting research flows at that rate carries substantially higher
    cost than the \gls{sciencedmz} approach and
    introduces latency, jitter, and state-tracking overhead detrimental to
    sustained bulk transfers. Router-based access control lists (\gls{acl}s)
    can provide an appropriate control profile for the well-characterised,
    high-throughput traffic profiles typical of research data transfers,
    when combined with dedicated monitoring and strict host
    hardening~\cite{dart2013}.
    \item[perfSONAR:] active network performance measurement
    platform~\cite{perfsonar}, enabling detection and diagnosis of
    problems.
    \item[Dedicated monitoring:] a specific \gls{ids} (Zeek, formerly Bro~\cite{zeek})
    monitors \gls{sciencedmz} traffic separately from the enterprise network.
\end{description}

This model is widely deployed across American research universities,
adopted at over 100 university campuses with funding from the NSF \emph{Campus
Cyberinfrastructure} (CC*) program~\cite{esnet-history,nsf-cc}. In Europe
and globally the same protected high-throughput function is met by
\gls{nren} dedicated circuits and lightpaths (\gls{geant}) and by
experiment networks (CERN LHCOPN/LHCONE), rather than by a US-style
campus enclave~\cite{geant,cern-lhcone}.

\smallskip\noindent\textbf{Limitations and compensating controls.}
The \gls{sciencedmz} model trades inline prevention for performance: Zeek
provides detection but not inline blocking, and \gls{acl}s do not inspect
application-layer content. These limitations are mitigated by (i)~\gls{soc}
monitoring of Zeek alerts, (ii)~endpoint hardening on \gls{dtn}s (minimal
services, regular patching), (iii)~the Trusted Activity Context
framework proposed in this report, which confines each research
activity to an isolated, monitored environment --- ensuring that a
compromise within one context cannot propagate laterally to other
activities or to the infrastructure core.

\subsubsection{Quantitative performance evidence}

\begin{table}[H]
\centering
\caption{Science DMZ case studies --- measured results}
\label{tab:sdmz-cases}
\small
\begin{tabularx}{\textwidth}{l X l}
\toprule
\textbf{University} & \textbf{Result} & \textbf{Ref.} \\
\midrule
U. of Florida & \textbf{x10} from the bandwidth upgrade alone,
    \textbf{x20} once the Science DMZ was added; 200~Gbps campus
    capacity for LHC/CMS data &
    \cite{esnet-uflorida} \\
\addlinespace
CU Boulder & \emph{Core Split} architecture separating research and
    admin from the border router; resolution of firewall
    bottlenecks & \cite{esnet-cuboulder} \\
\addlinespace
Penn State & From a 50~Mbps baseline, \textbf{$\sim$x5} (inbound) and
    \textbf{$\sim$x12} (outbound) gain after disabling the firewall
    TCP-flow-sequence-checking feature that violated RFC~1323
    (TCP Window Scaling) & \cite{esnet-pennstate} \\
\addlinespace
U. Michigan & Science DMZ with dedicated DTNs and perfSONAR; firewall
    repositioned near data sources rather than in the research
    path & \cite{umich-sdmz} \\
\addlinespace
Rutgers & Multi-campus fabric with Science DMZ + SDN + secure
    enclaves; research/central IT partnership &
    \cite{rutgers-sdmz} \\
\bottomrule
\end{tabularx}
\end{table}

Under classic TCP congestion control (Reno), the Mathis
equation~\cite{mathis1997} predicts that even minimal packet loss
severely degrades throughput on high-bandwidth links. While modern
algorithms such as CUBIC~--- which carries the majority of research
network traffic~--- are more loss-tolerant, they remain fundamentally
constrained by packet loss. Any loss or jitter introduced by stateful inspection
devices remains detrimental to sustained high-throughput transfers,
particularly on long-RTT paths~\cite{sciencedmz-perf}.

\subsubsection{Enterprise IT vs Science DMZ comparison}

\begin{table}[H]
\centering
\caption{Enterprise model vs Science DMZ comparison}
\label{tab:enterprise-vs-sdmz}
\small
\begin{tabularx}{\textwidth}{l X X}
\toprule
\textbf{Characteristic} & \textbf{Enterprise Model} &
    \textbf{Science DMZ} \\
\midrule
Philosophy & Prevention through blocking & Detection and rapid response \\
\addlinespace
Flow management & DPI of all packets & Lightweight \gls{acl}s + out-of-band \gls{ids} \\
\addlinespace
Performance impact & High (jitter, latency) & Minimal for intended bulk-transfer paths \\
\addlinespace
Tool management & Restrictive whitelist & Monitored autonomy \\
\addlinespace
Shadow IT risk & High & Reduced (fewer incentives for workarounds) \\
\bottomrule
\end{tabularx}
\end{table}

\subsection{VLAN network segmentation}
\label{subsec:vlan}

\gls{vlan} segmentation is standard practice in most major universities. It
consists of creating logically isolated network zones, each with its own
security policies.

\subsubsection{Typical VLAN categories}

\begin{table}[H]
\centering
\caption{VLAN categories in a typical university architecture}
\label{tab:vlans}
\small
\begin{tabularx}{\textwidth}{l l X l}
\toprule
\textbf{Zone} & \textbf{Security level} & \textbf{Content} & \textbf{Access} \\
\midrule
Administrative & High & Finance, HR, student IS, ERP &
    Restricted, \gls{ids}/\gls{ips} \\
\addlinespace
Research & Flexible & HPC, Science DMZ, lab networks &
    High throughput, direct NREN access \\
\addlinespace
Teaching & Moderate & Classrooms, LMS, teaching labs &
    Standard \\
\addlinespace
Staff BYOD & Intermediate & Staff personal laptops &
    Internet + progressive interop. \\
\addlinespace
Residential & Isolated & Student residences &
    NAT, limited, isolated \\
\addlinespace
IoT / Building & Restricted & HVAC, lighting, access control &
    Isolated, no Internet \\
\addlinespace
Guests / eduroam & Untrusted & Visitor Wi-Fi, conferences &
    Internet only \\
\bottomrule
\end{tabularx}
\end{table}

\smallskip\noindent\textbf{Important:} \gls{vlan}s provide \emph{logical}
separation but not security isolation on their own. They must be
combined with inter-\gls{vlan} filtering (\gls{acl}s or firewall rules at Layer~3
boundaries), DTP disabled on all access ports, the native \gls{vlan} changed
from default, and unused ports assigned to a dead \gls{vlan}. Without these
hardening measures, \gls{vlan} hopping attacks remain feasible. The Trusted
Activity Contexts proposed in this report add a further layer: even
within a \gls{vlan}, each activity operates in its own isolated, policy-bound
environment, limiting lateral movement \emph{within} a zone as well as
between zones.

\smallskip\noindent\textit{\footnotesize Technical implementation:
the network segmentation contract, zone catalogue, and Science DMZ
patterns are developed in Chapter~8 of the standalone Technical
Reference~\cite{ng-tech-ref-2026}. The wave-based migration approach
to declaring zones and operationalising the standard interfaces is
in Chapter~21 (Wave~3).}

\subsubsection{Multi-tier DMZ architecture}

Universities typically maintain several DMZ tiers:

\begin{itemize}[nosep]
    \item \textbf{External DMZ:} web servers, \gls{dns}, email relays ---
    directly exposed to the Internet;
    \item \textbf{Application DMZ:} application servers requiring
    limited external access;
    \item \textbf{\gls{sciencedmz}:} high-performance network path dedicated
    to research data transfers;
    \item \textbf{Management DMZ:} bastions (\emph{jump boxes}),
    monitoring systems, out-of-band management.
\end{itemize}

\subsection{Micro-segmentation and SDN}
\label{subsec:microseg}

Beyond traditional \gls{vlan} segmentation, modern approaches use:

\begin{description}
    \item[Micro-segmentation:] isolation at the individual workload
    level, and no longer only at the subnet level. Technologies:
    VMware NSX, Cisco ACI, Illumio.
    \item[SDN (\emph{Software-Defined Networking}):] enables dynamic
    provisioning of network segments and automated response to security
    incidents. Examples of platforms (commercial and open-source):
    Cisco DNA Center, OpenDaylight, ONOS.
    \item[Dynamic 802.1X:] port-level authentication with dynamic \gls{vlan}
    assignment based on the identity and posture of the device.
\end{description}

\subsection{Adapted Zero Trust architecture}
\label{subsec:zero-trust}

The \gls{zt} architecture (\emph{``never trust, always
verify''})~\cite{nist-zt} eliminates the concept of a trusted internal
network. Its application to universities requires specific
adaptations:

\begin{table}[H]
\centering
\caption{Progressive Zero Trust implementation in a university environment}
\label{tab:zero-trust}
\small
\begin{tabularx}{\textwidth}{l l X}
\toprule
\textbf{Phase} & \textbf{Component} & \textbf{Actions} \\
\midrule
1 & Identity & IAM, universal \gls{mfa}, federation via
    InCommon/eduGAIN~\cite{edugain} \\
\addlinespace
2 & Device & EDR, posture verification, device certificates \\
\addlinespace
3 & Micro-segmentation & SDN, least-privilege network access \\
\addlinespace
4 & Continuous monitoring & SIEM, UEBA, dynamic risk scoring \\
\bottomrule
\end{tabularx}
\end{table}

The key adaptation for universities is to recognise that:
\begin{itemize}[nosep]
    \item The requirement for managed devices is not achievable for all
    users (dominant \gls{byod});
    \item Differentiated access levels are necessary (more access with
    managed device + \gls{mfa}, limited access with unmanaged device);
    \item Research instruments and IoT require exemptions or special
    treatment;
    \item \gls{eduroam}~\cite{eduroam} and identity federations
    constitute the natural identity layer.
\end{itemize}

\subsection{Medical Science DMZ: proof for sensitive data}
\label{subsec:medical-dmz}

The argument ``everything must be locked down because of sensitive
data'' is challenged by the \keyword{Medical Science DMZ}, a design pattern
developed by researchers from Harvard, the University of Chicago,
Indiana University and Lawrence Berkeley National Laboratory~\cite{medical-sdmz}. This model handles
biomedical data (HIPAA) without sacrificing performance:

\begin{itemize}[nosep]
    \item Security through \textbf{end-to-end encryption} and strict
    endpoint authentication --- not through inspection of every packet;
    \item Sensitive data in \textbf{dedicated enclaves} with strict
    \gls{acl}s, logically separated from open research traffic.
\end{itemize}

If medical data subject to HIPAA can be processed in a segmented
architecture, university research data can \emph{a fortiori} be processed in such an architecture.

% ============================================================================
\section{Examples of universities with decentralised architecture}
\label{sec:exemples}
% ============================================================================

\epigraph{``If I have seen further, it is by standing on the shoulders of Giants.''}{--- Isaac Newton, letter to Robert Hooke, 1675/76~\cite{newton1675}}

This section presents concrete examples of world-leading universities
that successfully implement decentralised and federated architectures.

\medskip
\noindent\textbf{Relevant and Adequate Reference panel.} The panel
below covers institutions of varying size, governance models, and
regulatory environments. Its composition reflects four selection
criteria:
\begin{enumerate}[nosep]
    \item \textbf{Principles vs.\ implementations.} The architectural
    primitives discussed --- segmentation, enclave charters,
    federated identity, Science DMZ, BYOD zoning --- are governance
    and design choices that translate across scales of institution.
    \item \textbf{Shared backbone services.} All institutions in
    the panel access equivalent NREN backbone services (\gls{geant}
    in Europe, Internet2/ESnet in North America), interconnected
    into one global research fabric. The adopting institution's own
    NREN is a peer in that fabric (see Section~\ref{sec:lu-landscape}),
    providing equivalent primitives at the national scale.
    \item \textbf{Size-comparable institutions.} The panel includes
    \institution{EPFL} (comparable in size to a
    mid-sized research university)~\cite{epfl-rcp}, \institution{CERN}
    (governance model translatable to research-and-teaching
    consortia)~\cite{cern-network}, and Paris-Saclay
    (multi-institution federation in a public-research perimeter
    regulated under an EU national equivalent of the
    framework)~\cite{paris-saclay-cyber}, alongside larger
    institutions.
    \item \textbf{Track record.} Institutions of this kind commonly
    have substantial experience adapting external models to
    their specific public-research context.
\end{enumerate}

% --- MIT ---
\subsection{Massachusetts Institute of Technology (MIT)}

MIT operates one of the most historically decentralised IT models in
higher education. IS\&T (\emph{Information Systems \&
Technology})~\cite{mit-ist} is the central IT organisation, but MIT has
always adopted a philosophy of \emph{``distributed computing''} inherited
from the \keyword{Project Athena}
(1983--1991)~\cite{mit-athena}.

\begin{description}
    \item[Network:] \institution{MITnet} provides a high-throughput
    campus backbone connected to Internet2 and commercial
    ISPs~\cite{mit-facts}.
    \item[Departmental autonomy:] laboratories (CSAIL, Lincoln
    Laboratory, Media Lab, Koch Institute) operate their own IT staff,
    their own servers, and network segments. CSAIL's
    \emph{The Infrastructure Group}
    (TIG) independently manages multiple data
    centers~\cite{mit-csail-tig}.
    \item[Segmentation:] extensive use of \gls{vlan}s. Departments can request
    their own \gls{vlan}s; laboratories with sensitive data (ITAR, HIPAA)
    operate on isolated segments.
    \item[Security philosophy:] the MITnet Rules of
    Use~\cite{mit-network-rules} reflect a tradition of host-based
    security rather than perimeter defence. Each system owner is
    responsible for hardening their own systems.
    \item[Legacy:] MIT invented the \keyword{Kerberos} protocol during
    Athena~\cite{mit-kerberos}, still used on campus alongside
    Shibboleth/SAML.
\end{description}

% --- Stanford ---
\subsection{Stanford University}

Stanford operates a \keyword{federated} IT model. University IT
(UIT)~\cite{stanford-uit} provides the core infrastructure, but each of
the seven schools has its own IT organisation with significant autonomy.

\begin{description}
    \item[Network:] \institution{SUNet}~\cite{stanford-sunet} is the
    campus backbone, one of the most advanced university networks, with a
    multi-tier architecture (backbone, distribution, access).
    \item[Multiple IT organisations:]
    \begin{itemize}[nosep]
        \item The School of Engineering has its own IT group
        ($\sim$40~staff)~\cite{stanford-soe-it} managing \gls{hpc} clusters
        and laboratory networks;
        \item The School of Medicine operates Technology \& Digital
        Solutions (TDS)~\cite{stanford-som}, a parallel IT organisation
        with dedicated security staff, managing a HIPAA-compliant infrastructure;
        \item SLAC (\emph{National Accelerator Laboratory}) has an
        entirely independent network and IT~\cite{stanford-slac};
        \item The \institution{Stanford Research Computing Center}
        (SRCC)~\cite{stanford-srcc} provides centralised \gls{hpc} (Sherlock
        cluster) but individual laboratories also manage their own
        infrastructure.
    \end{itemize}
    \item[Security:] \emph{shared responsibility} model. UIT manages the
    perimeter firewall and campus policies, but each school's IT is
    responsible for security within its domain.
\end{description}

% --- Cambridge ---
\subsection{University of Cambridge}

Cambridge is one of the most \keyword{decentralised universities in the
world}, due to its collegiate
structure~\cite{cambridge-structure}.

\begin{description}
    \item[Network:] the \institution{UDN} (\emph{University Data
    Network}, historically called CUDN)~\cite{cambridge-udn} is the
    multi-gigabit fibre backbone connecting colleges, departments, and
    buildings. Connected to JANET/Jisc (UK \gls{nren}).
    \item[Collegiate decentralization:]
    \begin{itemize}[nosep]
        \item Each of the \textbf{31 colleges} operates its own IT
        infrastructure: network segments, servers, email systems, Wi-Fi;
        \item Departments (Computer Science and Technology, Cavendish
        Laboratory, MRC Laboratory of Molecular Biology) have substantial
        independent IT operations.
    \end{itemize}
    \item[Network delegation:] the UDN provides connectivity but
    \emph{delegates subnet management} to
    institutions~\cite{cambridge-connections}. Each institution manages
    its own \gls{dhcp}, \gls{dns}, and firewall rules.
    \item[IP delegation:] Cambridge owns a complete /16 block
    (131.111.0.0/16) and delegates /24 or larger blocks to
    institutions~\cite{cambridge-ip}.
    \item[Security:] the UDN operates as a largely open network within
    the backbone. Individual institutions are responsible for firewalling
    their own segments. UIS monitors the backbone for abnormal traffic
    and can isolate failing segments.
\end{description}

% --- ETH Zurich ---
\subsection{ETH Zurich}

ETH Zurich operates a \keyword{hybrid centralised-federated}
model~\cite{eth-departments}.
Informatikdienste (ITS) provides the central IT infrastructure, but each
of the 16 departments has its own IT Service Group
(ISG)~\cite{eth-isg-infk} and often dedicated IT staff.

\begin{description}
    \item[Network:] high-speed fibre backbone connecting the Zentrum and
    Hoenggerberg campuses, linked to SWITCH (Swiss \gls{nren}) at 100+ Gbps.
    \item[Autonomy:] each department can operate its own servers,
    including computing clusters. The Computer Science department (D-INFK)
    manages its own substantial server infrastructure.
    \item[Segmentation:] extensive \gls{vlan}s with 802.1X network access
    control. Separate zones: administrative, teaching, research, DMZ.
    \item[Federated services:] via SWITCH: SWITCHaai (federated identity),
    SWITCHdrive (cloud storage), SWITCHengines (OpenStack IaaS).
\end{description}

\textbf{Key innovation --- the Administrator Rights Agreement:}
Unlike a total ban, ETH Zurich has formalised root access through an
\emph{Agreement Regarding Administrator
Rights}~\cite{eth-admin-agreement}:
\begin{itemize}[nosep]
    \item The user can request root rights if needed for their thesis
    or research;
    \item By signing, the user accepts responsibility for maintenance;
    \item The user commits not to disable basic security mechanisms
    (logs, antivirus, updates);
    \item It is a \emph{transfer of responsibility} rather than a
    technical lockdown.
\end{itemize}

The baseline IT protection rules~\cite{eth-baseline} require that
non-compliant resources be placed in \emph{isolated network zones} so
that weaknesses are not exploitable from the outside.

% --- CMU ---
\subsection{Carnegie Mellon University (CMU)}

CMU has a historically significant role in distributed computing,
notably through the \keyword{Andrew Project} (1982--c.\,1989)~\cite{cmu-andrew}
which created the AFS distributed file system.

\begin{description}
    \item[Autonomy:] the School of Computer Science (SCS) operates its
    own IT infrastructure~\cite{cmu-scs-computing}, including its own
    Kerberos realm (\texttt{CS.CMU.EDU}), its AFS cell, its computing
    clusters, and its network segments.
    \item[SEI:] the Software Engineering Institute~\cite{cmu-sei},
    a DoD-funded Federally Funded Research and Development Center (FFRDC),
    operates an entirely separate IT environment with DoD security
    requirements.
    \item[Legacy:] AFS provides a global namespace
    (\texttt{/afs/cs.cmu.edu/})~\cite{cmu-scs-computing} with Kerberos
    authentication, historically used at many institutions (CERN used it
    for over 20~years before beginning migration to EOS/CERNBox).
\end{description}

% --- Michigan ---
\subsection{University of Michigan}

Michigan operates a \keyword{federated model with central
coordination}~\cite{umich-its, umich-it-community}.
ITS provides the shared infrastructure, but each of the 19 schools and
colleges has its own IT organisation.

\begin{description}
    \item[Network:] \institution{UMNet} provides the campus backbone
    across Ann Arbor, Dearborn and Flint.
    Michigan is also part of MiNet, a dedicated inter-university
    research network operated by Merit Network connecting U-M,
    Michigan State and Wayne State~\cite{umich-core-network}.
    Michigan co-founded Merit Network~\cite{umich-merit} and was among
    the founding universities of Internet2.
    \item[Federated governance:]
    \begin{itemize}[nosep]
        \item An \emph{IT Council}~\cite{umich-it-council} serves as an
        enterprise-wide advisory group to the VPIT-CIO, including IT
        leaders from each school;
        \item Michigan Medicine operates as a quasi-independent IT
        organisation (HITS) with its own
        CISO~\cite{umich-medicine-it};
        \item CAEN (\emph{Computer Aided Engineering
        Network})~\cite{umich-caen} in the School of Engineering has its
        own clusters and infrastructure;
        \item Advanced Research Computing
        (ARC)~\cite{umich-arc} provides centralised \gls{hpc} but laboratories
        can supplement with their own hardware.
    \end{itemize}
    \item[Micro-segmentation:] network divided into security zones with
    different trust levels. Sensitive data (CUI, HIPAA, FISMA) is placed
    on dedicated sensitive data
    environments~\cite{umich-safecomputing}.
    \item[Classification:] \emph{Safe Computing}~\cite{umich-safecomputing}
    program defining levels (High, Moderate, Low, Public) associated
    with specific network and system
    requirements.
\end{description}

% --- CERN ---
\subsection{CERN}

CERN operates one of the most sophisticated research network
architectures in the world~\cite{cern-network}.

\begin{description}
    \item[Campus network:] tens of thousands of connected devices across
    the Meyrin (Switzerland) and Prevessin (France) sites, with a
    multi-100~Gbps backbone and over 50\,000~km of optical fibre.
    \item[Experiment networks:] each LHC experiment (ATLAS, CMS, ALICE,
    LHCb) operates its own dedicated network segments:
    \begin{itemize}[nosep]
        \item Experiment control networks (isolated, high security);
        \item Data acquisition networks (ultra-high throughput, up to
        Tbps);
        \item Analysis networks (for physicists).
    \end{itemize}
    \item[Strict separation:] the \keyword{Technical Network} (accelerator
    control) is \emph{strongly isolated by firewalls and strict \gls{acl}s}
    from the general network (GPN) for critical system
    security~\cite{cern-tn}.
    \item[WLCG:] tiered architecture (Tier-0 at CERN, national Tier-1
    centres and university Tier-2/Tier-3 sites --- around
    160~computing centres in more than 40~countries) distributing
    worldwide computing via LHCONE~\cite{cern-lhcone, cern-wlcg}.
    \item[Security model:] CERN balances perimeter controls
    (closed-by-default outer firewall) with extensive \emph{monitoring
    and rapid response}, explicitly ensuring that security measures
    \emph{``do not obstruct the operations of CERN's accelerators and
    experiments or limit academic
    freedom''}~\cite{cern-security-freedom}.
\end{description}

% --- Oxford ---
\subsection{University of Oxford}

Oxford explicitly describes itself as \keyword{federated}~\cite{oxford-infosec}.
Departments have substantial autonomy and responsibility for implementing
information security under an overarching strategy defined by central
security functions. Unit heads are \emph{accountable}, may implement the
strategy as they see fit, and must define and document the
  roles, contacts and requirements of their department.

% --- EPFL ---
\subsection{École Polytechnique Fédérale de Lausanne (EPFL)}

EPFL has set up a \keyword{Research Computing Platform} (RCP)
offering an à la carte infrastructure~\cite{epfl-rcp}: storage,
GPU resources, research clusters and general server resources that
researchers can provision according to their needs.

\medskip
\noindent\textbf{ETH Zurich's Data Stewardship model:}
At ETH Zurich, \textbf{Data Stewards} are deployed directly in
departments~\cite{eth-data-stewardship} --- experts bridging the gap
between compliance and technical reality, ensuring security
\emph{by design} and not \emph{by obstruction}.

% --- Wisconsin-Madison ---
\subsection{University of Wisconsin-Madison}

The university promotes federation as a \emph{structured
partnership} preserving local research and teaching capabilities
while improving consistency and
monitoring~\cite{uw-madison}.

% --- NC State ---
\subsection{North Carolina State University}

Described by EDUCAUSE~\cite{ncsu-governance} as a federated IT
organisation with a central IT (OIT) and independent units in
colleges, libraries, institutes and departments. Governance is
explicitly aligned with the missions (education, research,
administration, user experience).

% --- Penn State ---
\subsection{Penn State}

Penn State has unified IT professionals from many units while
maintaining unit agency through participatory governance and
SLAs~\cite{pennstate-it}.

% --- French universities ---
\subsection{French universities}

\subsubsection{Université Paris-Saclay}

Paris-Saclay is a particularly interesting case because it is a
\keyword{grouping of institutions}~\cite{paris-saclay-network}: the
university proper, CentraleSupélec, ENS Paris-Saclay, IHES, and
numerous CNRS/CEA/INRIA laboratories.
\emph{Note: École Polytechnique, although geographically located on the
Saclay plateau, is a member of Institut Polytechnique de Paris (IP Paris),
a distinct entity from Université Paris-Saclay.}

\begin{itemize}[nosep]
    \item Each member institution maintains its own IT infrastructure
    and its own identity system;
    \item Federated identity uses the RENATER federation (FER) for
    inter-institutional authentication;
    \item The multi-institutional nature makes federation
    \emph{structural and not optional};
    \item Interoperability is ensured by standardised protocols
    (\gls{eduroam}, Shibboleth, RENATER connectivity).
\end{itemize}

\subsubsection{Sorbonne Université}

Resulting from the UPMC / Paris-Sorbonne merger, with integration of
associated research laboratories:

\begin{itemize}[nosep]
    \item The IT department provides central services, but faculty
    and laboratory IT teams persist;
    \item CNRS joint research units within Sorbonne ---
    such as LIP6~\cite{sorbonne-lip6} (computer science) and
    LPNHE~\cite{sorbonne-lpnhe} (nuclear physics) --- manage their own
    servers, computing clusters, and internal network configurations.
\end{itemize}

\subsubsection{CNRS--University model}

CNRS laboratories (UMR --- Joint Research Units), often
embedded within universities, \emph{naturally} create a federated
model: CNRS and university IT coexist, each UMR typically having
its own network segment within the host university's
infrastructure.

\subsection{Organizational models: charters and accountability}

Rather than informal exceptions, high-performing universities
use \keyword{institutional instruments}:

\begin{description}
    \item[Enclave charters:] modeled on Oxford~\cite{oxford-infosec},
    unit heads are accountable, define roles and contacts,
    and document their department's requirements.
    \item[Service catalogue:] EDUCAUSE~\cite{educause-catalog}
    provides a model including ``Research'', ``Teaching'' and
    ``Information Security'' as service categories.
    \item[Formalized exception path:] ETH Zurich formalizes
    self-assessment, registration, temporary exemptions and
    isolation zones as mechanisms to support non-standard
    needs~\cite{eth-baseline}.
    \item[IT charter:] in France, CNRS and universities
    such as UPPA~\cite{uppa-charte} or Le Havre~\cite{lehavre-charte}
    use charters defining responsibilities rather than
    blocking usage.
\end{description}

\textbf{Key persuasion point for central IT:} charters and
catalogs transform \emph{``freedom''} into \emph{accountable,
auditable autonomy}. They produce artifacts useful
as \gls{gdpr} Article~32 evidence.

% --- Summary table ---
\subsection{Comparative summary}

\begin{table}[H]
\centering
\caption{IT autonomy models in reference universities}
\label{tab:comparatif}
\small
\begin{tabularx}{\textwidth}{l c c X}
\toprule
\textbf{University} & \textbf{Model} & \textbf{Segm.} &
    \textbf{Distinctive characteristic} \\
\midrule
MIT & Distributed & VLAN & Athena philosophy, host-based security \\
Stanford & Federated & Multi-tier & 7 schools with autonomous IT, independent SLAC \\
Cambridge & Decentralised & Delegated & 31 autonomous colleges, delegated IP \\
ETH Zurich & Hybrid & 802.1X & Per-department ISG, federated SWITCH \\
CMU & Distributed & VLAN & Global AFS, autonomous SCS \\
Michigan & Coordinated fed. & Micro-seg. & IT Council, 19 schools \\
CERN & Multi-domain & Isolated (FW) & Per-experiment networks, global WLCG \\
Oxford & Federated & Departmental & Accountable unit heads, charters \\
EPFL & Hybrid & Enclaves & RCP, Data Stewards in labs \\
Wisconsin-Madison & Federated & Partnership & Federation preserving local capabilities \\
NC State & Federated & Missions & Mission-aligned governance \\
Penn State & Unified-federated & SLA & Participatory governance \\
Paris-Saclay & Structural fed. & Institutional & Multi-institution federation \\
\bottomrule
\end{tabularx}
\end{table}

% ============================================================================
\section{Media-Covered Failures in Universities (2024--2025): Context and Lessons}
\label{sec:defaillances-mediatiques}
% ============================================================================

\epigraph{``The first principle is that you must not fool yourself~--- and you are the easiest person to fool.''}{--- Richard P.\ Feynman, \textit{Cargo Cult Science}, 1974~\cite{feynman1974}}

Proposals for \keyword{network segmentation}, \keyword{federated governance}
and differentiated access (research, \gls{byod}, AI) are sometimes opposed on the
grounds that ``open'' or ``flexible'' policies would expose universities to
the kind of high-profile incidents that made headlines in 2024 and 2025.
\textbf{This section summarises the most media-covered failures, analyses
their root causes, and shows that these incidents do not support the claim
that the solutions proposed in this report are too risky} --- on the
contrary, they underscore the need for \emph{segmentation} and
\emph{risk-based governance} rather than for uniform restriction.

\subsection{Scope and purpose}

The aim here is not to minimise the harm caused by breaches or to assign
blame to specific institutions. It is to provide a factual basis so that
debate relies on \textbf{cause and architecture}, not on anecdote. When
incidents are cited as proof that ``permissive'' policies fail, the
appropriate response is to distinguish: (i)~\emph{third-party and
supply-chain vulnerabilities}, (ii)~\emph{centralised, unsegmented
perimeters} enabling lateral movement, and (iii)~\emph{credential
compromise} --- none of which are addressed by locking down research and
teaching zones. Sector reports (Verizon DBIR, Sophos, EDUCAUSE) and
normative frameworks (NIST) consistently point to \keyword{segmentation}
and \keyword{isolation of activity contexts} as defences, not to blanket restriction~\cite{verizon-dbir2024,sophos-ransomware-edu2025,educause-isg,nist-csf}.

\subsection{Most media-covered incidents (2024--2025)}

\subsubsection{France}

\begin{description}
    \item[\institution{Université Paris-Saclay} (August 2024):] A
    ransomware attack claimed by RansomHouse affected central
    administrative services, email and data servers; the institution's
    website and digital services were severely disrupted. The attack
    reportedly involved theft of around 1\,TB of data and compromise of
    employee credentials (infostealer activity). Recovery lasted months; the
    university worked with the French national
    cybersecurity agency (ANSSI). \textbf{Relevant point:} the systems
    targeted were \emph{central} administrative and email infrastructure ---
    a large, undifferentiated perimeter. The incident illustrates the
    risk of a \keyword{single point of failure} and the value of
    isolating sensitive administrative data in a dedicated enclave, as
    recommended in this report~\cite{paris-saclay-cyber,soccer-paris-saclay}.
    \item[\institution{Université de Rennes} (March 2025):] The university
    reported a cyberattack affecting a \emph{teaching subnet} (ISTIC ---
    computing and electronics). The group FunkSec claimed theft of
    sensitive data and used double extortion. The university stated that
    the impact was limited to that subnet and that main networks remained
    operational. \textbf{Relevant point:} containment to a subnet is
    consistent with the benefit of \keyword{network segmentation}; the
    incident does not support the argument that flexible or segmented
    policies are inherently riskier~\cite{rennes-cyber}.
    \item[\institution{Universit\'e Paris~1 Panth\'eon-Sorbonne}:]
    Reported an October~2024 cyberattack in which directory data of
    about 73\,000 students and staff (names, email addresses, photos,
    status) was extracted; the university notified the national
    data-protection authority~\cite{paris-sorbonne-cyber}.
\end{description}

\subsubsection{United States and international (2021--2024)}

\begin{itemize}[nosep]
    \item \textbf{Texas Tech University Health Sciences Center}
    (September 2024): Ransomware (Interlock); about 1.4 million individuals
    affected; healthcare and administrative
    data~\cite{ttuhsc-breach2024}. One of the largest
    university-related breaches of the year; the context is a \emph{health
    sciences} and administrative perimeter, not research openness.
    \item \textbf{MOVEit supply-chain attacks} (2023--2024): The Cl0p
    group exploited a vulnerability in Progress MOVEit Transfer software.
    Many universities were affected (e.g.\ UCLA, Stanford Medicine and the
    University System of Georgia, among many others). \textbf{These breaches were caused by a \emph{third-party
    software vulnerability}, not by institutional policies on openness or
    segmentation.} Affected institutions patched the product and improved
    vendor oversight; the lesson is supply-chain and vendor risk, not
    restriction of research or teaching zones~\cite{moveit-breach,verizon-dbir2024}.
    \item \textbf{Howard University} (September 2021): Ransomware forced
    a campus shutdown and class cancellations for approximately ten days;
    the university stated that personal data did not appear to have been
    compromised~\cite{howard-breach2021}. \textbf{Western Sydney
    University} (2024): Multiple incidents; threat actors used
    \emph{compromised IT accounts} to access the student management system
    and data warehouse~\cite{wsu-breach2024}. This pattern --- \textbf{credential compromise and
    lateral movement across systems} --- is exactly what \keyword{network
    segmentation} is intended to limit~\cite{elisity-edu}.
\end{itemize}

\subsubsection{2025: Oracle E-Business Suite and other vectors}

In 2025, a major wave of incidents in higher education was linked to
exploitation of \emph{critical zero-day vulnerabilities in Oracle
E-Business Suite} (CVSS~9.8), used by the Cl0p group
for mass data extortion~\cite{oracle-ebs-edu2025}. Affected institutions included the University
of Phoenix (3.5~million records, per its SEC filing), Harvard
(1.3~TB exfiltrated), Dartmouth (35\,000+ individuals), University of
Pennsylvania, and Baker University, among
others~\cite{oracle-ebs-edu2025}. Oracle EBS is widely used for
\emph{financial and administrative operations} (payroll, tuition, vendor
payments). \textbf{Again, the root
cause is a \emph{third-party product vulnerability} and supply-chain risk,
not the degree of openness of research or teaching networks.} Sector
surveys (e.g.\ Sophos State of Ransomware in Education 2025) report that
in higher education, \emph{exploited vulnerabilities} and \emph{unknown
security gaps} are among the top identified causes --- consistent with
the need for hardening, patching, and segmentation of high-value
perimeters, not with uniform lockdown~\cite{oracle-ebs-edu2025,sophos-ransomware-edu2025}.

\subsection{Root causes: what the failures actually show}

\begin{table}[H]
\centering
\caption{Recurring causes of media-covered university incidents (2024--2025)}
\label{tab:causes-incidents}
\small
\begin{tabularx}{\textwidth}{>{\hsize=0.8\hsize}X >{\hsize=1.2\hsize}X}
\toprule
\textbf{Cause type} & \textbf{What it implies for policy} \\
\midrule
Third-party / supply-chain (MOVEit, Oracle EBS) &
    Vendor risk, patching, and isolation of critical apps --- not
    restricting research/teaching zones. \\
\addlinespace
Credential compromise \& lateral movement (e.g.\ Paris-Saclay, Western Sydney) &
    Need for \keyword{segmentation} and least-privilege access so that
    one compromise does not expose the whole institution. \\
\addlinespace
Centralised, undifferentiated perimeter (e.g.\ Paris-Saclay central services) &
    Aligns with this report's argument: a single large perimeter
    creates a \keyword{single point of failure}; administrative data
    should be in a dedicated, hardened enclave. \\
\addlinespace
Phishing / human factor (a recurring vector in higher education~\cite{verizon-dbir2024}) &
    \gls{mfa}, awareness, and identity governance --- not blanket
    restriction of academic networks. \\
\bottomrule
\end{tabularx}
\end{table}

The Verizon 2024 Data Breach Investigations Report and NIST guidance
emphasise \textbf{segmentation and isolation} to limit lateral
movement and contain incidents~\cite{verizon-dbir2024,nist-sp800-215}.
\textbf{Media-covered failures in 2024--2025 do not demonstrate that
segmented, risk-based architectures are too risky; they demonstrate that
\emph{lack} of segmentation and \emph{undifferentiated} perimeters
increase impact when a breach occurs.}

\begin{tcolorbox}[
    colback=OliveGreen!5!white,
    colframe=OliveGreen!75!black,
    title=\textbf{Key takeaway},
    breakable
]
\textbf{The incidents most cited in the media (Paris-Saclay, MOVEit,
Oracle EBS, Texas Tech, Rennes, etc.) are not evidence that the
solutions proposed in this report --- segmentation, federated
governance, differentiated zones --- ``do not work'' or are ``too
risky''.} Taken together, these incidents more strongly support the
need for segmentation, supply-chain risk management, credential
protection, and internal containment than for uniform restriction.
The root causes --- third-party and supply-chain risk, credential
compromise, and centralised perimeters without internal boundaries
--- require \emph{precisely the kind of architecture and governance}
advocated here: strong
protection of the administrative enclave, clear boundaries between
zones, and accountable autonomy in research and teaching, with
monitoring and incident response aligned to each zone's risk level.
\end{tcolorbox}

\subsection{Conclusion of this section}

A fair reading of the most media-covered university failures in 2024 and
2025 supports \textbf{strengthening segmentation and federated
governance}, not retreating into uniform restriction. Opponents of
differentiated policies may invoke these incidents as proof of risk;
this section provides the factual and analytical basis to show that
\textbf{the causes are largely unrelated to ``open'' research or
teaching policies} and that the recommended measures (multi-zone
architecture, charters, evidence-based controls) are the appropriate
response to the types of failures observed.

% ============================================================================
\section{National Research and Education Networks (NREN)}
\label{sec:nren}
% ============================================================================

\epigraph{``The web is more a social creation than a technical one. I designed it for a social effect~--- to help people work together~--- and not as a technical toy.''}{--- Tim Berners-Lee, \textit{Weaving the Web}~\cite{bernerslee1999}}

Universities are not isolated: they belong to research and
education network ecosystems that provide connectivity,
federated services and shared security.

\subsection{European ecosystem}

\begin{description}
    \item[GÉANT~\cite{geant}:] pan-European backbone network connecting
    around 40 \gls{nren}s across Europe, serving 50 million users
    in 10\,000+ institutions. Backbone capacity: up to 500~Gbps.
    \item[\Loc{nren-name}\loc{cite-restena}:] the adopter's own national \gls{nren}.
    \item[RENATER~\cite{renater}:] French \gls{nren} connecting all
    universities, research laboratories (CNRS, INRIA, CEA, INSERM)
    and French grandes écoles. Multi-100~Gbps backbone.
    \item[DFN~\cite{dfn}:] German \gls{nren} with the X-WiN backbone and DFN-CERT for
    incident response.
    \item[SWITCH~\cite{switch-nren}:] Swiss \gls{nren} with SWITCHaai (federated identity),
    SWITCHdrive and SWITCHengines (OpenStack IaaS).
    \item[JANET/Jisc:] UK \gls{nren}.
\end{description}

\subsection{Key federated services}

\begin{table}[H]
\centering
\caption{Federated services provided by NRENs}
\label{tab:nren-services}
\small
\begin{tabularx}{\textwidth}{l l X}
\toprule
\textbf{Service} & \textbf{Operator} & \textbf{Description} \\
\midrule
eduroam~\cite{eduroam} & GÉANT & Federated Wi-Fi: a researcher
    connects at any worldwide institution with their local credentials \\
\addlinespace
eduGAIN~\cite{edugain} & GÉANT & Global identity interfederation
    (80+ federations, inter-institutional SSO) \\
\addlinespace
FER & RENATER & French identity federation (part of eduGAIN) \\
\addlinespace
GÉANT TCS & GÉANT & Trusted certificate service for
    institutions \\
\addlinespace
perfSONAR~\cite{perfsonar} & ESnet/GÉANT & Distributed network
    performance monitoring \\
\addlinespace
Globus~\cite{globus} & U. Chicago & High-performance research data
    transfer \\
\bottomrule
\end{tabularx}
\end{table}

\subsection{NREN architecture: common pattern}

All \gls{nren}s follow a similar architectural pattern:

\begin{enumerate}[nosep]
    \item \textbf{National backbone:} high-speed optical network
    connecting major cities;
    \item \textbf{Regional aggregation:} regional networks or points of
    presence (PoP);
    \item \textbf{Campus connection:} each university connected at
    10--100~Gbps;
    \item \textbf{Peering:} interconnection between \gls{nren}s (via \gls{geant} in
    Europe) and with the commercial Internet;
    \item \textbf{Federated services:} identity, \gls{eduroam}, certificates,
    security (\gls{csirt}).
\end{enumerate}

\textbf{Essential point:} this model inherently relies on
\keyword{federation} and \keyword{decentralization}. If \gls{nren}s
operate on a federated model at the national and international level,
it is natural and consistent for universities to adopt the same model
internally.

\subsection{NREN status: what is typically in place, what is missing}
\label{subsec:lux-sdmz-status}

\loc{sdmz-status-lead} public documentation of \loc{nren-name} and
of \loc{institution-hpc-facility} does not, as of early 2026, evidence
a deployed \gls{sciencedmz} at \loc{institution-name} or at the national
\gls{nren} level. A typical NREN advertises 10--100~Gbps
point-to-point links on demand and a 100~Gbps uplink to
\gls{geant}~\cite{geant}, but publishes no dedicated \gls{dtn} fabric,
no public \keyword{perfSONAR} endpoint~\cite{perfsonar}, and no
Science DMZ reference design. HPC data-transfer documentation commonly
lists \texttt{rsync}, \texttt{scp}, \texttt{sshfs} and SMB to central
storage --- a profile typical of an enterprise-IT path, not the
friction-free, monitored research path described
in~\cite{dart2013, esnet-cases}.

\medskip
\noindent\textbf{The substrate is present; the pattern is not.}
The \emph{building blocks} required for a \gls{sciencedmz} are
frequently already operational: high-bandwidth \gls{geant}
connectivity through \loc{nren-name}, \loc{nren-csirt} for
incident response, and \loc{institution-hpc-facility} with
petabyte-scale research traffic to \loc{national-hpc-facility}
and to international partners.
What is missing is the architectural \emph{pattern}: a dedicated
zone, \gls{dtn}s, perfSONAR observability, and the governance
charter that together turn raw bandwidth into research-usable,
monitored, segregated capacity.

\medskip
\noindent\textbf{Is it a reasonable option?} On the
evidence available, yes --- with two qualifications that the
design must internalise:

\begin{itemize}[nosep]
    \item \textbf{Mutualised, not siloed.} Where the national
    research ecosystem is small (the institution together with
    \loc{research-ecosystem-partners}), it may be too small for
    several separate Science DMZs. The operationally rational
    form is then an \gls{nren}-anchored Science DMZ \emph{co-designed}
    with the national NREN, federated across the research-active
    institutions, with a Medical Science DMZ
    enclave~\cite{medical-sdmz} for biomedical workloads
    subject to \gls{gdpr} Article~9.
    \item \textbf{Benefits are not automatic.} Recent
    quantitative work shows that a Science DMZ outperforms
    enterprise paths only after deliberate tuning of hosts,
    protocol stacks and measurement; absent that tuning,
    latency on long-RTT paths can even
    degrade~\cite{sciencedmz-perf}. The investment must be
    paired with operational ownership (perfSONAR baselining,
    Zeek monitoring, \gls{dtn} lifecycle), not treated as a
    one-shot procurement.
\end{itemize}

The operative question is therefore not ``Should the
institution adopt a Science DMZ on its own?'' but: \emph{Should
the national ecosystem, through its \gls{nren}, deploy a federated
research-network zone that the institution and the other
research-active institutions can jointly consume?}
This reframing is consistent with the federation principle just
established and with the cost-feasibility structure developed in
Section~\ref{sec:couts}.

% ============================================================================
\section{Normative and governance frameworks}
\label{sec:normes}
% ============================================================================

\epigraph{``Plans are worthless, but planning is everything.''}{--- Dwight D.\ Eisenhower, 1957~\cite{eisenhower1957}}

\subsection{NIST Cybersecurity Framework 2.0}

The \framework{NIST Cybersecurity Framework}~\cite{nist-csf} (version 2.0,
February 2024) provides a structured approach to cybersecurity risk
management, widely adopted in higher education.

\begin{table}[H]
\centering
\caption{NIST CSF 2.0 functions applied to universities}
\label{tab:nist-csf}
\small
\begin{tabularx}{\textwidth}{l l X}
\toprule
\textbf{Function} & \textbf{Code} & \textbf{University application} \\
\midrule
Govern & GV & Digital Policy Commission with elected academic
    representation, risk appetite acknowledging research openness \\
\addlinespace
Identify & ID & Asset inventory (BYOD challenge), data classification,
    risk assessment by unit \\
\addlinespace
Protect & PR & IAM via InCommon/Shibboleth, \gls{mfa}, training, technologies
    that do not impede research workflows \\
\addlinespace
Detect & DE & SOC (shareable via REN-ISAC~\cite{ren-isac}),
    Zeek/Suricata monitoring, threat intelligence \\
\addlinespace
Respond & RS & Response plans accounting for the academic calendar
    and research continuity \\
\addlinespace
Recover & RC & Continuity plans for critical periods
    (exams, enrollment), research data recovery \\
\bottomrule
\end{tabularx}
\end{table}

\textbf{Key point:} the NIST CSF is a \emph{risk}-based framework,
not a \emph{restriction}-based one. It does not prescribe uniformly
reducing the attack surface through prohibitions --- rather, it prescribes
assessing risks and applying proportionate controls.

\subsection{ISO 27001:2022}

The \framework{ISO/IEC 27001:2022}~\cite{iso27001} standard provides a
framework for an Information Security Management System (ISMS).

\subsubsection{University adaptations}

\begin{itemize}
    \item \textbf{Scope (Clause 4):} the \gls{isms} can cover central IT
    only or specific systems, reflecting the reality
    of federated governance;
    \item \textbf{Leadership (Clause 5):} management commitment
    must include the president/rector and the vice-president for teaching,
    not only the CIO;
    \item \textbf{Control A.8.22 (Network segregation):} directly
    relevant to the \gls{vlan} / \gls{sciencedmz} segmentation described above;
    \item \textbf{Control A.8.23 (Web filtering):} particularly
    sensitive in an academic context because web filtering may conflict
    with academic freedom.
\end{itemize}

\subsubsection{ISO 27001 certified universities}

Several universities --- including the University of Leeds and the
University of Warwick (UK, IT services) and Eindhoven University of
Technology (Netherlands), among others in Germany, Australia and
Japan --- have obtained certification, typically for specific scopes.

\textbf{Observation:} in most cases, certification covers
\emph{central IT services} and not the entire institution, which
\emph{reflects and validates} the federated governance model.

\subsection{NIST SP 800-171 and secure enclaves}

For universities receiving sensitive funding (notably
US DoD), \framework{NIST SP 800-171} Rev.\,3~\cite{nist-171r3} (May 2024) specifies
97 security requirements (consolidated from 110 in Rev.\,2~\cite{nist-171}). The university response is \emph{not}
to lock everything down, but to create compliant \keyword{enclaves}
for processing controlled data, while the rest of
the institution operates at security levels appropriate to each zone.

This \textbf{secure enclave} model within an open architecture
has accelerated the adoption of network segmentation and \gls{zt}
in research universities.

\subsection{COBIT 2019 for IT governance}

The \framework{COBIT 2019}~\cite{cobit} framework from ISACA is
particularly relevant for universities because it explicitly
distinguishes \emph{governance} (strategic direction) from
\emph{management} (operational execution).

\begin{table}[H]
\centering
\caption{IT decision domains adapted to universities (after Weill
\& Ross~\cite{weill-ross})}
\label{tab:cobit}
\small
\begin{tabularx}{\textwidth}{l l X}
\toprule
\textbf{Domain} & \textbf{Who decides} & \textbf{Example} \\
\midrule
IT principles & Leadership + Digital Policy Commission & Cloud-first policy,
    open-source preference \\
\addlinespace
IT architecture & Central IT (with faculty input) & Network standards,
    identity federation \\
\addlinespace
IT infrastructure & Central IT & Data centers, backbone, baseline security \\
\addlinespace
Application needs & Faculties/departments & Research software, specialised
    databases \\
\addlinespace
IT investments & Joint committee & Budget allocation between research IT
    and administrative IT \\
\bottomrule
\end{tabularx}
\end{table}

\subsection{A nationally provided risk-analysis methodology}

Several jurisdictions provide an open-source, \framework{ISO/IEC
27005}-aligned risk-analysis method and tool --- often developed by
\loc{national-cyber-competence-centre} --- that enables organisations
to conduct optimised, precise and repeatable risk assessments,
leveraging the capitalisation of risk-analysis scenarios already
produced in similar contexts, and integrated into the national
incident-notification platform as the risk-analysis module for
NIS2-related obligations.

For the adopting institution, \loc{national-risk-method}\loc{cite-national-risk-method} carries a specific
operational relevance: being aligned with \framework{ISO/IEC 27005},
it is directly usable for the institutional maturity-report structure
and the per-zone risk-analysis evidence that Proposition~14
(Section~\ref{sec:propositions}) consolidates into a single annual
score submitted to the national supervisory authority. Where no such
national tool exists, any ISO/IEC~27005-conformant method serves the
same purpose.

\subsection{Risk-based approach vs.\ compliance-based approach}

\begin{table}[H]
\centering
\caption{Comparison of security approaches}
\label{tab:risque-conformite}
\small
\begin{tabularx}{\textwidth}{l X X}
\toprule
\textbf{Criterion} & \textbf{Compliance/restriction approach}
    & \textbf{Risk-based approach} \\
\midrule
Driver & Regulatory checklists & Threat and asset assessment \\
\addlinespace
Key question & ``Are we meeting the controls?'' &
    ``What are our critical assets and likely threats?'' \\
\addlinespace
Drawback & ``Checkbox'' mentality, uniform &
    Requires risk analysis skills \\
\addlinespace
Benefit for univ. & Low & \textbf{High} --- enables
    differentiation by zone \\
\addlinespace
Associated frameworks & Internal control lists & NIST CSF, ISO 27001,
    FAIR risk quantification~\cite{fair} \\
\bottomrule
\end{tabularx}
\end{table}

\textbf{Proposition:} the optimal approach combines both ---
regulatory compliance (\gls{gdpr}, sector-specific regulations) as a
\emph{floor}, and risk assessment to prioritise
investments \emph{above} that floor according to institutional
risk appetite.

% ============================================================================
\section{Artificial intelligence and access to cloud services}
\label{sec:ia}
% ============================================================================

\epigraph{``We can only see a short distance ahead, but we can see plenty there that needs to be done.''}{--- Alan Turing, \textit{Computing Machinery and Intelligence}~\cite{turing1950}}

\subsection{AI as a structural need for research and teaching}

In 2025-2026, generative artificial intelligence and foundation
models have become everyday tools for:

\begin{itemize}[nosep]
    \item \textbf{Research}: training and fine-tuning models,
    accessing provider APIs (OpenAI GPT, Anthropic Claude,
    Google Gemini, Meta LLaMA, Mistral), using cloud platforms
    (AWS SageMaker, Google Vertex AI, Azure ML, HuggingFace Hub);
    \item \textbf{Teaching}: machine learning and NLP courses
    requiring API access, practicals on foundation models,
    comparative evaluation of AI systems;
    \item \textbf{Cybersecurity research}: AI-based threat analysis,
    anomaly detection, adversarial testing.
\end{itemize}

\subsection{The GDPR challenge and non-EU cloud services}

Use of major cloud-based AI services (OpenAI, Anthropic, Google
Cloud AI) may involve international transfers of personal data and
therefore requires case-by-case assessment of data flows, transfer
safeguards, and technical and organisational measures appropriate
to the risk. For research, this creates a \textbf{structural tension}
between:
\begin{itemize}[nosep]
    \item The obligation to protect personal data (\gls{gdpr});
    \item The need for free access to the best global AI platforms
    to remain scientifically competitive.
\end{itemize}

\subsection{Architectural solution: differentiated AI access zones}

\begin{table}[H]
\centering
\caption{AI access policy by network zone}
\label{tab:ia-zones}
\small
\begin{tabularx}{\textwidth}{l X X}
\toprule
\textbf{Zone} & \textbf{Authorized AI access} &
    \textbf{Legal framework} \\
\midrule
Administrative & Validated institutional AI only (e.g.: Copilot
    with university tenant) & Strict \gls{gdpr}, no personal data
    to non-EU services \\
\addlinespace
Teaching & AI APIs for courses (with anonymized or
    synthetic data) & Instructor agreement, no real student data \\
\addlinespace
Research & \textbf{Free access} to major cloud-based AI services
    & Researcher-centred accountability per
    Section~\ref{sec:safeguards-ai}; cumulative \emph{Schrems~II}
    regime applies to personal-data transfers~\cite{schrems2020} \\
\addlinespace
BYOD & Via the zone corresponding to the authenticated profile &
    According to the destination zone \\
\bottomrule
\end{tabularx}
\end{table}

\subsection{Safeguards for research AI access}
\label{sec:safeguards-ai}

\textbf{Academic-community-centred design within the regulatory
framework.} The researcher-centred accountability model presented
below operationalises \gls{gdpr} Articles~24(1) and~5(2) and the
AI Act's research-support principle (Recital~25)~\cite{eu-ai-act}.
Six items detail the framework:

\begin{enumerate}[nosep]
    \item \textbf{Researcher accountability principle.} The Principal
    Investigator (PI) is responsible for the compliance of their
    project with applicable \gls{gdpr},
    \emph{Schrems~II}~\cite{schrems2020} and AI
    Act~\cite{eu-ai-act} provisions. The institution provides
    classification templates, registration facilities and training;
    the PI applies scientific and legal judgment case by case.

    \item \textbf{Data classification.} Before any transfer to a
    non-EU AI cloud provider, the PI classifies the data:
    \begin{itemize}[nosep]
        \item \emph{Non-personal or fully anonymised data}:
        free transfer (\gls{gdpr} not applicable);
        \item \emph{Pseudonymised personal data}: transfer permitted
        under \gls{scc} + documented \gls{tia} + supplementary
        measures (encryption with key separation, reinforced
        contractual clauses);
        \item \emph{Identifying personal data}: transfer requires
        explicit informed consent (or another Art.~6 lawful basis)
        + \gls{scc} + supplementary measures + documented
        \gls{tia}, or the data must remain within the EEA.
    \end{itemize}

    \item \textbf{Cooperative declaration via institutional register.}
    On first use of a non-EU AI cloud provider for personal data, the
    central IT service registers the access on PI request, recording:
    provider, data classification (provided by the PI), lawful basis
    (declared by the PI), and safeguards in place. The register is
    maintained as a standing institutional artefact available for
    periodic \gls{dpo} audit. The model is \textbf{declarative and
    cooperative}, not gatekeeping: the PI requests access, the IT
    service registers, the DPO audits.

    \item \textbf{EU-first routing preference.} European providers
    (e.g.\ Mistral, OVH, Scaleway) are preferred where functionality
    is equivalent. International collaboration with non-EU providers
    is permitted as an \emph{academic right} under the declarative
    model above.

    \item \textbf{Environment separation, traceability, training.}
    Research AI notebooks and environments operate within the
    research zone, isolated from the administrative zone; external
    API access is logged for audit; PIs receive training on
    \gls{gdpr}, \emph{Schrems~II}~\cite{schrems2020}, the AI
    Act~\cite{eu-ai-act}, and the data-classification framework.

    \item \textbf{AI Act application.} The risk-based classification
    of the AI Act~\cite{eu-ai-act} applies; many research uses fall
    outside high-risk categories. Bespoke research tools developed
    for scientific R\&D are exempt under Article~2(6); research,
    testing and development before market placement is exempt under
    Article~2(8). The report endorses Recital~25 calling for research
    support and the avoidance of undue restriction.
\end{enumerate}

\medskip
\noindent\textit{Advocacy positioning.} The cooperative declaration
model proposed in item~3 operationalises the \gls{gdpr} accountability
principle in a research context where strict per-access gatekeeping
would unduly restrict international scientific collaboration. This
report indicates that the institution could share its register
template and audit methodology with the national supervisory
authority, \loc{national-dpa} and EDPB working
groups, as a contribution to
research-context regulatory practice (see
Section~\ref{sec:advocacy}, Position~D.5).

% ============================================================================
\section{Proposed infrastructure and organisational solutions}
\label{sec:propositions-infra}
% ============================================================================

\epigraph{``In anything at all, perfection is finally attained not when there is no longer anything to add, but when there is no longer anything to take away.''}{--- Antoine de Saint-Exup\'ery, \textit{Terre des Hommes}, 1939~\cite{saintexupery1939}}

Building on the state of the art and the previous examples, this
section proposes a target architecture for a university wishing to
reconcile security and academic freedom.

\subsection{Conceptual architecture: five structural layers}
\label{sec:concept-arch}

Before detailing the operational zones, it is useful to present the
\textbf{conceptual model} that underpins them. Following established
resilience-engineering principles~\cite{avizienis2004,nist-sp800-160v2},
a university information system can be decomposed into five structural
layers, each with distinct dependability requirements:

\begin{enumerate}
    \item \textbf{Infrastructure core (dependability foundation).}
    The set of services on which \emph{all other layers depend}: identity
    and access management (\gls{iam}), \gls{dns}, \gls{dhcp}, \gls{pki}, NTP,
    \textbf{immutable backup systems} (air-gapped or append-only
    repositories --- the primary defence against ransomware),
    central storage, hypervisors, orchestration platforms,
    \textbf{\gls{siem}} for centralised security event aggregation
    and alerting (e.g.\ Wazuh, OpenSearch),
    network access control (\gls{nac}),
    campus backbone, Wi-Fi controllers, \gls{vpn} concentrators,
    and administration bastions. This core constitutes the
    \emph{dependability foundation} of the entire system: it must be
    strongly segmented, hardened, redundant, supervised,
    and isolated from experimental networks. Its compromise would affect
    every other layer; its resilience guarantees the survival of the
    institution~\cite{nist-sp800-160v2,nist-csf}.

    \item \textbf{Common services.}
    Shared platforms consumed by the entire university community:
    messaging, \gls{lms}, user storage, collaborative tools, videoconferencing,
    printing, intranet portals, \gls{hpc} scheduling interfaces. These services
    must be highly available, redundant, backed up, and organised in
    Trusted Activity Contexts aligned with data sensitivity --- but they
    do not require the same isolation as the infrastructure core.

    \item \textbf{Autonomous zones.}
    Perimeters managed locally by research laboratories, departments,
    experimental platforms, scientific clusters, fablabs, or temporary
    project infrastructures. Each zone enjoys \emph{administrative
    autonomy} within a dedicated network segment, under a responsibility
    charter (ETH~Zurich model~\cite{eth-admin-agreement}). The guiding
    principle is: \emph{local autonomy, bounded global impact}.

    \item \textbf{Exposed zones.}
    Public-facing services (web servers, open platforms, external
    collaborative portals, APIs, partner gateways) that must be
    \emph{isolated from the internal SI} by DMZ-type
    architectures~\cite{nist-sp800-53}.

    \item \textbf{User networks.}
    End-user connectivity: staff wired, student wired, \gls{eduroam} Wi-Fi,
    guest Wi-Fi, \gls{byod}, administrative workstations, computer labs. Open
    by nature, but segmented from critical zones.
\end{enumerate}

\noindent
The seven operational security zones detailed below map onto these five
conceptual layers. The administrative zone and sensitive-data enclave sit
within the infrastructure core and common-services layers; the research
and teaching zones are autonomous zones; the \gls{sciencedmz} and \gls{byod} zone
span the exposed and user-network layers respectively.

\subsection{Trusted Activity Contexts}
\label{sec:propositions-tac}

\begin{tcolorbox}[
    colback=NavyBlue!4!white,
    colframe=NavyBlue!75!black,
    boxrule=0.6pt,
    sharp corners=south,
    title={\textbf{The transverse role of Trusted Activity Contexts}}
]
The five layers of Section~\ref{sec:concept-arch} are \emph{spatial}:
they describe \emph{where} components live and \emph{how strongly}
they are isolated from one another. The \keyword{Trusted Activity
Context} (\gls{tac}) introduced in Section~\ref{sec:tac} is by
contrast \emph{transverse}: it describes \emph{what coherent,
validated environment} an activity operates in, regardless of which
layer hosts that activity. A research project, a teaching session,
or an administrative process can each be a \gls{tac} that consumes
resources from several layers --- identity and \gls{pki} from the
infrastructure core, collaborative tools from the common services,
autonomous compute from a research zone --- with policy gating at
every step and audit correlating the whole. The operational
realisation of \gls{tac}s --- end-to-end activation sequence, policy
contract, endpoint posture, conformance suite --- is developed in
the companion Technical Reference~\cite{ng-tech-ref-2026}.
\end{tcolorbox}

\subsection{Segmented network architecture with BYOD zone and AI access}

\begin{figure}[H]
\centering
\begin{tikzpicture}[
    node distance=0.7cm and 1.2cm,
    zone/.style={rectangle, draw, rounded corners, minimum width=2.6cm,
                 minimum height=0.9cm, font=\scriptsize, align=center},
    core/.style={rectangle, draw, rounded corners, minimum width=4cm,
                 minimum height=0.9cm, font=\small\bfseries, align=center,
                 fill=gray!20},
    inet/.style={ellipse, draw, minimum width=4.5cm, fill=blue!10,
                 font=\small},
    arr/.style={-{Stealth[length=2mm]}, thick},
    darr/.style={{Stealth[length=2mm]}-{Stealth[length=2mm]}, thick}
]
    \node[inet] (inet) {Internet / NREN / GÉANT / AI Clouds};
    \node[core, below=0.8cm of inet] (border) {Border router};
    \node[core, below=0.8cm of border] (core) {Campus network core};

    % Row 1 - 5 zones (horizontal offsets reduced to narrow the figure)
    \node[zone, below left=1cm and 2.6cm of core, fill=red!15]
        (admin) {Admin\\(High\\security)};
    \node[zone, below left=1cm and -0.5cm of core, fill=orange!15]
        (enseign) {Teaching\\(Standard)};
    \node[zone, below=1cm of core, fill=green!15]
        (rech) {Research\\(Flexible)};
    \node[zone, below right=1cm and -0.5cm of core, fill=cyan!15]
        (sdmz) {Science\\DMZ};
    \node[zone, below right=1cm and 2.6cm of core, fill=purple!15]
        (byod) {BYOD\\Personal};

    % Row 2 - details
    \node[zone, below=0.6cm of admin, fill=red!5]
        (admin-d) {ERP, HR\\Strict GDPR};
    \node[zone, below=0.6cm of enseign, fill=orange!5]
        (enseign-d) {LMS, VDI\\Cyber-range};
    \node[zone, below=0.6cm of rech, fill=green!5]
        (rech-d) {HPC, GPU\\Free AI APIs\\Autonomous labs};
    \node[zone, below=0.6cm of sdmz, fill=cyan!5]
        (sdmz-d) {DTN\\perfSONAR};
    \node[zone, below=0.6cm of byod, fill=purple!5]
        (byod-d) {Internet\\+ progressive\\interop.};

    % Row 3 - extra
    \node[zone, below=0.6cm of admin-d, fill=yellow!10]
        (iot) {IoT/Bldg.\\(Isolated)};
    \node[zone, below=0.6cm of rech-d, fill=Maroon!10]
        (enclave) {Sensitive\\data\\enclave};

    % Arrows
    \draw[darr] (inet) -- (border);
    \draw[darr] (border) -- (core);
    \draw[arr] (core) -- (admin);
    \draw[arr] (core) -- (enseign);
    \draw[arr] (core) -- (rech);
    \draw[arr] (core) -- (sdmz);
    \draw[arr] (core) -- (byod);
    \draw[arr] (admin) -- (admin-d);
    \draw[arr] (enseign) -- (enseign-d);
    \draw[arr] (rech) -- (rech-d);
    \draw[arr] (sdmz) -- (sdmz-d);
    \draw[arr] (byod) -- (byod-d);

\end{tikzpicture}
\caption{Segmented network architecture with BYOD zone and AI access. The
Science~DMZ realises the generic high-throughput research-data path; an
\gls{nren}-operated dedicated circuit or \gls{geant} lightpath is an
equivalent realisation.}
\label{fig:archi-proposee}
\end{figure}

\subsection{Seven security zones}

\begin{enumerate}
    \item \textbf{Administrative zone (high security):}
    \begin{itemize}[nosep]
        \item ERP, HR, finance, student information systems;
        \item Stateful firewall, \gls{ids}/\gls{ips}, mandatory proxy
        (note: HTTPS inspection on the mandatory proxy must comply with
        the \emph{B\u{a}rbulescu}~\cite{barbulescu2017} proportionality criteria regarding
        staff personal communications);
        \item Strict \gls{gdpr} compliance, database activity monitoring,
        encrypted storage at rest;
        \item Restricted access, mandatory \gls{mfa};
        \item \emph{This is the zone that must be the most locked down ---
        not the entire university.}
    \end{itemize}

    \item \textbf{Teaching zone (standard security):}
    \begin{itemize}[nosep]
        \item \gls{lms} (\emph{Learning Management System}), classrooms,
        teaching laboratories;
        \item \gls{vdi} (\emph{Virtual Desktop Infrastructure}) for specialised
        courses --- technologies: VMware Horizon, Citrix, or
        Apache Guacamole as HTML5 access gateway over
        Proxmox~VE (lightweight open-source alternative to
        commercial \gls{vdi});
        \item Containerization (Docker, Kubernetes) for practical
        environments: JupyterHub on Kubernetes~\cite{z2jh} allows providing
        each student with an isolated environment for data science,
        statistics, and computational courses across all faculties;
        \item Cyber-ranges for security and networking courses, as
        recommended by NIST~\cite{nist-cyberrange};
        \item Access to AI APIs with anonymized data --- available to
        all disciplines (AI-assisted legal research, computational
        linguistics, data analysis in social sciences, etc.);
        \item Standard security controls, \gls{eduroam} Wi-Fi;
        \item Instructors across all faculties must be able to freely
        configure teaching environments.
    \end{itemize}

    \item \textbf{Research zone (flexible, high throughput):}
    \begin{itemize}[nosep]
        \item \gls{hpc} clusters (Slurm job scheduler, parallel filesystem
        such as Lustre or BeeGFS, high-speed interconnect),
        GPU servers, laboratory infrastructure;
        \item \textbf{Administrative autonomy} for laboratories and
        departments: each unit manages its own servers, software
        and configurations within its dedicated \gls{vlan} (central
        network operations retain control of \gls{vlan} boundaries and
        inter-\gls{vlan} routing);
        \item Apptainer (formerly Singularity)~\cite{singularity} for \gls{hpc}
        containers without root privileges;
        \item \textbf{Free access to major cloud-based AI services}
        (see Section~\ref{sec:ia});
        \item Monitoring-based security (Zeek + Suricata) rather than
        restriction-based, with log export to the central \gls{siem};
        \item Direct \gls{nren} access for massive data transfers.
    \end{itemize}

    \item \textbf{High-throughput research-data path --- \gls{sciencedmz} instantiation (high-performance transfers):}
    \begin{itemize}[nosep]
        \item \gls{dtn} with Globus for inter-institutional transfers,
        requiring Linux kernel TCP tuning (buffer sizes, IRQ affinity,
        jumbo frames) to achieve line-rate on 100~Gbps links;
        \item \gls{acl}-based security (no stateful firewall, incompatible with
        10--100~Gbps flows) on dedicated switching infrastructure;
        \item perfSONAR for performance monitoring (nodes placed both
        inside and outside the \gls{sciencedmz} to identify boundary
        bottlenecks);
        \item Zeek + Suricata for dedicated surveillance (protocol
        analysis and signature-based detection in parallel).
        \item[] \emph{The Science~DMZ is one instantiation of this role;
        where a campus enclave is not warranted, an \gls{nren}-operated
        dedicated circuit or \gls{geant} lightpath realises the same
        high-throughput research-data path.}
    \end{itemize}

    \item \textbf{Personal \gls{byod} Zone:}
    \begin{itemize}[nosep]
        \item Full Internet access;
        \item Progressive interoperability with other zones depending on
        the authenticated profile (see Section~\ref{sec:specificites}),
        implemented via FreeRADIUS and 802.1X dynamic \gls{vlan} assignment;
        \item \gls{mfa} + posture verification for extended access:
        for unmanaged devices, posture is established through device
        certificates or MDM enrolment; agentless checks (OS fingerprinting
        via DHCP, browser-based detection) provide a baseline but cannot
        verify patch level or antivirus status;
        \item Network access control (PacketFence or equivalent) manages
        device onboarding, \gls{vlan} placement, and policy enforcement.
    \end{itemize}

    \item \textbf{IoT/Building Zone (isolated):}
    \begin{itemize}[nosep]
        \item HVAC, lighting, physical access control;
        \item Strictly isolated from the campus IP network
        (dedicated \gls{vlan} with no default gateway, or physically
        separate switching for access-control systems);
        \item Local NTP and firmware-update server within the zone
        (devices cannot reach the Internet but must still be patched
        and time-synchronised);
        \item Single-point gateway (application-layer proxy) for
        facilities-management access to building controllers.
    \end{itemize}

    \item \textbf{Sensitive data enclave:}
    \begin{itemize}[nosep]
        \item For projects requiring enhanced compliance
        (medical data, classified projects, etc.);
        \item Micro-segmentation, enhanced logging, strictly
        controlled access;
        \item Compliant with NIST SP 800-171 if necessary;
        \item \emph{Enclave within the open architecture} --- the model
        proven by universities receiving DoD funding and by
        Harvard's Medical \gls{sciencedmz}.
    \end{itemize}
\end{enumerate}

\begin{principlebox}[NIS2 holistic scope and Article~21 proportionality]
All seven zones described above operate \emph{within} the \gls{nis2}
perimeter applicable to the institution as an essential entity
(Section~\ref{sec:nis2}). The variation in control profiles between
zones is the direct expression of \gls{nis2} Article~21(1), which
requires measures \emph{appropriate and proportionate} to each
entity's exposure to risk. National transpositions typically make
this proportionality principle binding and extend the calibration
criteria to include the state of the art, applicable European and
international standards, cost of implementation, entity size, and
the likelihood and severity of incidents. The seven-zone segmentation
is the operational mechanism through which the institution implements
its NIS2 obligations across a heterogeneous public-research perimeter.

\smallskip
\textbf{Segmentation does not narrow the regulatory scope; it
calibrates the implementation.} Every zone is subject to risk
management, incident reporting, supply-chain controls and continuous
evidence production --- the calibration is in the intensity of
controls, not in the inclusion within scope.
\end{principlebox}

\subsection{Federated IT Governance}
\label{subsec:gouvernance-federee}

\medskip
\noindent\textbf{Continuous-improvement mechanism.} The Digital
Policy Commission (Section~\ref{sec:dpc}), \loc{academic-governing-bodies}, the responsibility charters and the ICRADP
(Section~\ref{sec:commission-ajustements}) collectively constitute
the institution's \keyword{continuous-improvement mechanism}. The
institutional artefact through which this mechanism becomes
opposable is the annual compliance and maturity report submitted
to the national supervisory authority pursuant to the national
NIS2 transposition, which closes each annual
governance cycle.

\subsubsection{Proposed Structure}

\begin{table}[H]
\centering
\caption{Distribution of responsibilities in a federated model}
\label{tab:gouv-federee}
\small
\begin{tabularx}{\textwidth}{l X X}
\toprule
\textbf{Domain} & \textbf{Central IT} &
    \textbf{Research units / Departments / Laboratories} \\
\midrule
Backbone network & Operates and maintains & Uses \\
\addlinespace
Identity (IAM) & Provides the infrastructure & Manages local accounts \\
\addlinespace
Baseline security & Defines minimum standards, operates the SOC &
    Implements within its perimeter \\
\addlinespace
Research servers & Provides optional hosting &
    \textbf{Freely administers} within its VLAN \\
\addlinespace
Research software & --- & \textbf{Freely installs} \\
\addlinespace
ERP / HR / Finance & Operates and secures & --- \\
\addlinespace
HPC & Provides the central cluster & Can supplement with own hardware \\
\addlinespace
Research cloud & Provides IaaS (OpenStack) &
    Provisions its own VMs \\
\addlinespace
BYOD zone & Provides the network infrastructure & Uses \\
\addlinespace
AI cloud & Maintains the institutional register (per PI request)
    and operates the periodic DPO audit & PI classifies data,
    requests access, uses the services under the cooperative
    declaration model (Section~\ref{sec:safeguards-ai}) \\
\addlinespace
\gls{gdpr} compliance & Coordinates, DPO & Responsible for its data \\
\bottomrule
\end{tabularx}
\end{table}

\subsubsection{Executive training and management responsibility}
\label{sec:exec-training}

The national NIS2 transposition imposes specific obligations on
the management bodies of essential entities, distinct from
researcher-level training and from the institutional security
architecture:
\begin{itemize}[nosep]
    \item \textbf{Executive training} (\gls{nis2}
    Art.~20(2)~\cite{nis2}): \loc{executive-leadership} and the
    members of \loc{governing-board} must receive cybersecurity
    training enabling informed governance decisions on the matters
    covered by the NIS2 obligation framework. The training is a
    governance enablement obligation, not an operational compliance
    burden.
    \item \textbf{Personal liability of management bodies}
    (\gls{nis2} Art.~20~\cite{nis2}): under the national
    transposition, members of the management body
    can be held personally accountable for failures to discharge
    the cybersecurity-risk management obligations.
\end{itemize}

\noindent The Digital Policy Commission proposed in
Section~\ref{sec:dpc} is the natural delivery mechanism for the
executive-training programme: its mixed academic-technical
composition allows the training to be structured as a series of
governance briefings tailored to the institution's specific
public-research perimeter, rather than a generic compliance course.
The training programme is documented in the maturity report
submitted to the national supervisory authority, providing
auditable evidence of compliance with NIS2 Art.~20(2).

\noindent\textit{Distinction.} The training of researchers on
\gls{gdpr}, \emph{Schrems~II} and the AI Act
(Section~\ref{sec:safeguards-ai}, item~5) is a separate and
parallel programme; the present sub-subsection covers
governance-body training only.

\subsubsection{Digital Policy Commission}
\label{sec:dpc}

The governance of digital infrastructure and AI policy should be
entrusted to a \keyword{Digital Policy Commission} whose key
characteristics are:
\begin{itemize}[nosep]
    \item \textbf{Chair:} an elected academic --- not a member of
    the IT operational hierarchy;
    \item \textbf{Composition:} a compact body of representatives
    from the departments and interdisciplinary centres across the
    university's faculties, ensuring that research, teaching, and
    interdisciplinary perspectives are all represented;
    \item \textbf{Advisory members (non-voting):} the CIO, the CISO,
    the \gls{dpo}, and a student representative;
    \item \textbf{Decision path:} the Commission formulates policy
    proposals that are submitted to \loc{academic-endorsement-path}
    for endorsement before implementation.
\end{itemize}

This Commission holds \emph{decision-making authority} over digital
and AI policy proposals that affect academic missions. The CIO/CISO
submit proposals and provide technical expertise; the Commission
\emph{decides} whether those measures are proportionate and compatible
with academic needs.

\smallskip\noindent
This model rests on the principle of \keyword{university democracy}:
the UNESCO \emph{Recommendation concerning the Status of
Higher-Education Teaching Personnel} (1997)~\cite{unesco1997}
explicitly affirms the right of academic staff to \emph{``elect a
majority of representatives to academic bodies within the higher
education institution.''}  The AAUP \emph{Statement on Government of
Colleges and Universities} (1966)~\cite{aaup1966} establishes that in
matters affecting teaching and research, the faculty should exercise
\emph{``primary responsibility,''} with governing board review only
\emph{``in exceptional circumstances.''} If elected academic bodies
have decision-making authority over curricula, research priorities, and
hiring, it is coherent that they also govern the digital infrastructure
that conditions all of these activities.

\smallskip\noindent
In practice, governance models vary: Oxford's Information and Digital
Committee~\cite{oxford-infosec} reports to Council with genuine
decision-making authority over digital strategy; Michigan's IT
Council~\cite{umich-it-council} and Harvard~FAS's IT Governance
Committee~\cite{harvard-it-governance} operate in a consultative
capacity. The University of Kentucky's Senate has created an AI
governance committee with policy-making authority derived from the
elected Senate~\cite{uk-senate-ai}. In California, elected faculty
senates hold statutory \emph{primacy} over technology-related academic
matters~\cite{asccc-ai}. This separation of strategic direction
(governance) from operational execution (management) mirrors
COBIT~2019's foundational distinction~\cite{cobit}. The Commission
also defines the minimum standards that each zone must comply with.

\smallskip\noindent
Applied to \loc{institution-name} --- \loc{institution-faculty-structure} with strong industrial and institutional
partnerships --- the proposed Commission would include representatives
from these diverse entities. Its proposals would follow the
institutional decision path: endorsement by \loc{academic-endorsement-path} before operational
implementation by the IT services.

\subsubsection{Contextual role management and separation of privileges}
\label{subsubsec:contextual-roles}

In a university, a single individual routinely holds \emph{multiple
roles}: a doctoral student may also be a teaching assistant, a part-time
system administrator for the laboratory, and a collaborator on an
external project. The classical approach of assigning one account with
the \emph{union} of all privileges violates the
\keyword{principle of least privilege} and makes audit trails
opaque~\cite{sandhu1996}.

The foundational \gls{rbac} model~\cite{sandhu1996} introduces the concept of
\keyword{session-based role activation}: a user does not permanently
carry all assigned roles but \emph{activates a subset per session},
depending on the task at hand. Applied to a university context, this
means:

\begin{itemize}[nosep]
    \item A researcher accessing the \gls{hpc} cluster activates her
    \emph{researcher} role with access to the research zone and AI
    cloud APIs;
    \item The same person, when entering grades, activates her
    \emph{instructor} role with access to the \gls{lms} in the teaching zone;
    \item When performing system administration on her laboratory's
    servers, she activates her \emph{local administrator} role, scoped
    to her unit's \gls{vlan};
    \item At no point does she hold simultaneous access to
    administrative ERP data, which requires a separate role governed by
    a distinct Trusted Activity Context.
\end{itemize}

\noindent
This approach yields several benefits aligned with the proposed
architecture:

\begin{itemize}[nosep]
    \item \textbf{Least privilege:} each session carries only the
    permissions necessary for the current task;
    \item \textbf{Traceability:} audit logs record \emph{which role}
    was active during each action, not merely which person;
    \item \textbf{Separation of duties:} conflicting roles (e.g.\
    ``data subject'' and ``data processor'') can be declared mutually
    exclusive within a single session~\cite{sandhu1996};
    \item \textbf{\gls{tac} alignment:} role activation maps naturally onto the
    Trusted Activity Context structure --- each role grants access to a
    specific context within the zone architecture.
\end{itemize}

\noindent
Examples of commonly deployed \gls{iam} platforms that support the
access-control dimension of this model include: Keycloak
(session-scoped role activation via OIDC/SAML), and Microsoft
Entra~ID (Privileged Identity Management for just-in-time role
elevation). Shibboleth serves as the SAML federation bridge to
\gls{edugain}~\cite{edugain} and releases attributes per service
provider, but does not itself manage session-scoped role selection. The full \gls{tac} vision ---
composing VLAN assignment, storage, applications, and access rights
from a single role activation event --- requires an orchestration
layer (portal or \gls{api} gateway) that coordinates these platforms;
no single turnkey product achieves this today, but the individual
building blocks are mature and deployable. The federated identity
layer (\gls{edugain} for SAML-based \gls{sso}, \gls{eduroam} for
network authentication) provides the inter-institutional foundation.
For modern applications and cloud APIs, OIDC/OAuth2 (natively
supported by Keycloak) complements the SAML layer.

\smallskip\noindent
\textbf{Note on Active Directory environments:} where Windows AD is
used, group membership is static per session. Implementing contextual
role activation in AD-joined environments requires either Microsoft
Entra~ID Privileged Identity Management for role elevation, or
dedicated access-brokering solutions. The \gls{tac} model applies
most naturally to web-based and Linux-based research environments,
and can be progressively extended to AD-managed workstations as
the identity infrastructure matures. A more detailed (and deliberately
tentative) endpoint integration matrix across Linux, macOS, and
Windows is sketched in Chapter~2 of the standalone Technical
Reference~\cite{ng-tech-ref-2026}; the broader endpoint contract
and the iOS/Android coverage gap are addressed in Chapter~11.

\subsubsection{Decision flow for autonomy requests}

A predictable process replaces informal exceptions:

\begin{enumerate}[nosep]
    \item \textbf{Qualification:} does the request concern
    administration, teaching, research, or \gls{byod}?
    \item \textbf{Risk assessment:} sensitive data? External exposure
    required? Cloud access outside the EU?
    \item \textbf{Zone assignment:} the request is directed to
    the appropriate zone with its native controls.
    \item \textbf{Charter:} if local autonomy is necessary
    (root admin, servers, AI APIs), a charter formalizes the
    responsibilities, incident contacts, and the monitoring baseline.
    \item \textbf{Validation:} the Digital Policy Commission validates
    cases involving external exposure or sensitive data, with
    endorsement by \loc{academic-endorsement-path} for
    institution-wide implications.
\end{enumerate}

\subsection{Cloud and Virtualization for Research}

\begin{description}
    \item[OpenStack private cloud:] allows researchers to provision
    virtual machines on demand, with isolation per tenant.
    Examples: Jetstream2 (NSF, Indiana University), Nectar
    (Australia), EGI Federated Cloud (Europe), SWITCHengines (Switzerland).

    \textbf{Note:} the recommendation is primarily to \emph{consume}
    existing research cloud infrastructure provided by \gls{nren}s
    (SWITCHengines, EGI Federated Cloud, \gls{geant} services) or to deploy
    lighter-weight alternatives (OpenNebula, CloudStack) for local
    needs. Operating a full OpenStack deployment requires significant
    dedicated operations capacity and is best suited to consortium-level
    or \gls{nren}-hosted deployments rather than individual university
    installations.

    \item[Kubernetes:] container orchestration for:
    \begin{itemize}[nosep]
        \item JupyterHub: interactive computing environments for
        students and researchers;
        \item Kubeflow: ML/AI pipelines;
        \item BinderHub: reproducible environments from
        Git repositories.
    \end{itemize}
    \item[VDI for teaching:] Apache Guacamole (open-source,
    clientless HTML5 access), VMware Horizon, or Citrix to provide
    standardised desktops accessible from any student
    device (\gls{byod}).
\end{description}

\subsection{Shared Security Organizations}

Rather than managing everything alone, a university can rely on:

\begin{itemize}[nosep]
    \item \textbf{REN-ISAC} (US / North America)~\cite{ren-isac}: threat intelligence
    sharing, DNS RPZ feeds, incident notifications (700+
    member institutions);
    \item \textbf{Trusted CI}~\cite{trustedci}: center of excellence
    in cybersecurity for research (funded by the NSF);
    \item \textbf{GEANT Security}: security services for
    European institutions;
    \item \textbf{the national / governmental \gls{csirt}}: the
    public-sector incident-response team; the institution's
    primary first-line CSIRT as a public institution;
    \item \textbf{the national NREN's CSIRT}: the \gls{csirt}
    of the national \gls{nren}, commonly a \emph{Trusted Introducer}
    member, complementary to \loc{national-gov-csirt} for incidents
    arising on or affecting the academic network perimeter;
    \item \textbf{a private-sector CSIRT}, where one exists, serving
    the private sector and non-governmental entities;
    relevant when third-party private partners of the
    institution are involved (not the institution's first-line
    CSIRT);
    \item \textbf{FIRST}~\cite{first}: global forum of incident
    response teams (\gls{csirt}), including DFN-CERT, SWITCH-CERT
    and national NREN and governmental CSIRTs;
    \item \textbf{JRES}~\cite{jres}: biennial French conferences on
    academic networks, a major source of experience feedback.
\end{itemize}

% ============================================================================
\section{Justification and Feasibility}
\label{sec:couts}
% ============================================================================

A university entering the third millennium will be, above all, a
\textbf{digital university}. The quality, security, and flexibility of
its digital infrastructure condition every academic activity: research,
teaching, administration, and international collaboration. For \loc{staffing-profile}, it
is both coherent and necessary that this standard of excellence extend
to the digital infrastructure that underpins all academic missions.

\medskip\noindent
The investment required is not a cost to be minimised but a
\textbf{strategic enabler}: just as a university invests in libraries,
laboratories, and lecture halls, it must invest in a digital
infrastructure worthy of its ambitions. The progressive approach
proposed in this dossier (governance first, then pilot on selected
entities, then gradual extension) is designed to be implementable
within the institution's \loc{funding-instrument},
spreading effort over several cycles and aligning each step with
demonstrated results.

\subsection{Cost analysis: centralised approach}

\begin{itemize}
    \item \textbf{Direct costs:} oversized centralised infrastructure
    having to handle all use cases, expensive application firewall licenses
    processing all traffic including research;
    \item \textbf{Indirect costs:}
    \begin{itemize}[nosep]
        \item Loss of researcher productivity (time spent
        circumventing restrictions or waiting for authorizations);
        \item Inability to teach certain advanced courses;
        \item Brain drain to more flexible universities;
        \item Shadow IT creating unmanaged risks;
        \item Loss of competitiveness for grant applications
        requiring high-performance research infrastructure.
    \end{itemize}
\end{itemize}

\subsection{Cost analysis: federated approach}

\begin{itemize}
    \item \textbf{Enhanced security:} Trusted Activity Contexts limit the
    propagation of incidents --- a compromise in a research context
    cannot reach the administrative context;
    \item \textbf{Proportionate controls:} maximum investment where
    the risk is highest (sensitive administrative data),
    adapted investment elsewhere;
    \item \textbf{Research productivity:} researchers work
    without hindrance in their zone, while being protected;
    \item \textbf{Pedagogical excellence:} advanced courses have
    realistic environments;
    \item \textbf{Attractiveness:} the university attracts high-level
    researchers thanks to a modern and flexible infrastructure;
    \item \textbf{AI access:} research and teaching at world-class level;
    \item \textbf{Functional \gls{byod}:} staff productivity with
    personal devices;
    \item \textbf{Facilitated compliance:} secure enclaves
    allow meeting the strictest requirements without
    penalizing the rest of the institution;
    \item \textbf{\gls{gdpr}/\gls{nis2} evidence:} each zone produces its own
    evidence of compliance and effectiveness.
\end{itemize}

\subsection{Feasibility and cost structure of the proposed approach}

The question of \textbf{cost and feasibility} is addressed explicitly,
in line with \gls{gdpr} Article~32, which requires that measures be
\emph{appropriate to the risk} taking into account \emph{state of the
art}, \textbf{costs} and context~\cite{gdpr}.

\begin{itemize}[nosep]
    \item \textbf{Proportionate investment:} by concentrating the
    strictest controls on the \emph{administrative enclave} only, the
    university avoids over-investing in perimeter security for
    low-sensitivity research and teaching traffic; the ``high security''
    perimeter becomes \emph{smaller and cheaper to harden and monitor}.
    \item \textbf{Reuse of existing ecosystems:} many components are
    \emph{low-cost or cost-free}: connectivity and security services
    from the NREN and its continental research backbone (in Europe,
    \gls{geant}), REN-ISAC and \gls{geant} Security (threat intelligence,
    \gls{csirt}), open-source tools (perfSONAR, Zeek, Apache Guacamole,
    JupyterHub, Globus). \gls{sciencedmz} patterns are documented and
    adopted by 100+ universities~\cite{esnet-history,nsf-cc} without requiring proprietary
    enterprise licenses for research zones.
    \item \textbf{Progressive implementation:} starting with governance
    (Digital Policy Commission, charters, service catalogue), followed
    by a pilot deployment on the entities with the strongest domain
    expertise (\loc{cs-research-unit}, \loc{security-research-centre}, and the institution's HPC team),
    then gradual extension to other faculties, departments and centres
    as results are validated (see Section~\ref{sec:transition}).
    \item \textbf{Comparative cost factors:} indirect costs
    associated with centralised approaches (productivity loss,
    talent mobility, Shadow IT, competitiveness in calls for
    proposals) and incident-related costs (remediation, fines,
    reputation impact in case of a major breach) provide a
    benchmark against which the proposed investment can be
    assessed. The analysis of media-covered university incidents
    (Section~\ref{sec:defaillances-mediatiques}) identifies the
    absence of internal boundaries as a recurring cause.
\end{itemize}

A detailed transition plan with phases, risk register and KPIs
(Section~\ref{sec:transition}) allows the institution to \emph{control
scope and budget} at each step and to prioritise quick wins
(governance, charters, pilot) before broader rollout.

% ============================================================================
\section{Transition Plan}
\label{sec:transition}
% ============================================================================

\noindent
\emph{An indicative scoping of a minimum-viable \gls{tac} pilot ---
with the caveat that timelines and resourcing remain to be confirmed
by the IT department --- is developed in the standalone Technical
Reference~\cite{ng-tech-ref-2026}, Chapter~23 (``Wave~5 --- TAC
pilot''). The full wave-based migration playbook spanning ten waves
is developed in Part~III of the same document.}

\subsection{Progressive Transition}

\begin{table}[H]
\centering
\caption{Progressive transition plan}
\label{tab:transition}
\small
\begin{tabularx}{\textwidth}{l X}
\toprule
\textbf{Step} & \textbf{Actions} \\
\midrule
\textbf{0. Initial declarations} &
    Auto-registration with the national supervisory authority
    (national NIS2 transposition) \\
& Initial setup of the institutional account on
    the national notification platform (incident-notification + risk-analysis modules) \\
& Definition of the maturity-report template
    (an ISO/IEC~27005-aligned risk-analysis method, such as a
    nationally provided open-source tool) and
    identification of per-zone evidence sources \\
\midrule
\textbf{1. Governance} &
    Establishment of the Digital Policy Commission
    (elected chair + faculty/centre representatives) \\
& Definition of target zones and charter templates \\
& Adoption of the differentiated responsibility charter
    (ETH Zurich model) \\
& Service catalogue and formalised exception path \\
& Executive-training programme for \loc{governing-board}
    (\gls{nis2} Art.~20(2), see Section~\ref{sec:exec-training}) \\
\midrule
\textbf{2. Pilot} &
    Pilot deployment of the segmented architecture and Trusted Activity
    Contexts on the entities with the strongest domain expertise (\loc{cs-research-unit}, \loc{security-research-centre}, and the institution's HPC team) \\
& Installation of \gls{dtn} + perfSONAR + Zeek \\
& PoC micro-segmentation (\gls{sdn} overlay) \\
& Pilot \gls{byod} zone with progressive interoperability \\
& Before/after metrology for factual data \\
\midrule
\textbf{3. Progressive extension} &
    Gradual rollout to other departments, faculties and
    interdisciplinary centres \\
& Establishment of local ISGs (IT correspondents trained by the
    IT services) \\
& Progressive \gls{zt} deployment (\gls{mfa} $\rightarrow$ posture
    $\rightarrow$ micro-segmentation) \\
& AI zone with worldwide cloud access for research \\
& Audit and continuous improvement cycle \\
\midrule
\textbf{4. Annual cycle} &
    Annual maturity-report submission to the national supervisory
    authority via the national notification platform (per the
    national NIS2 transposition) \\
& Internal-audit cycle covering all enclaves, feeding the
    maturity report \\
& \gls{dpo} audit cycle covering the cooperative declaration
    register for non-EU AI cloud access
    (Section~\ref{sec:safeguards-ai}) \\
& Executive-training refresh for \loc{governing-board}
    (\gls{nis2} Art.~20(2), Section~\ref{sec:exec-training}) \\
\bottomrule
\end{tabularx}
\end{table}

\noindent The governance step is predominantly organisational and can be
carried out with existing staff. The pilot confines initial deployment
to the entities with the strongest domain expertise --- \loc{cs-research-unit}, a cybersecurity-focused research centre,
and the institution's HPC team --- whose expertise in
computer science, cybersecurity, and high-performance infrastructure
makes them natural candidates for validating the architecture.
Progressive extension then follows validated results, allowing the
institution to control scope at each step.

\subsection{Risk Register}

\begin{table}[H]
\centering
\caption{Transition programme risk register}
\label{tab:risques}
\small
\begin{tabularx}{\textwidth}{X l l X}
\toprule
\textbf{Risk} & \textbf{Prob.} & \textbf{Impact} &
    \textbf{Mitigation} \\
\midrule
Central IT perceives enclaves as loss of control &
    High & High & Show the security gain on the admin enclave;
    charters and precedents (Oxford, ETH) \\
\addlinespace
Departments treat autonomy as absence of responsibility &
    Medium & High & Charters with accountability, mandatory contacts,
    baseline monitoring, periodic reviews \\
\addlinespace
Exceptions creating a ``hidden flatness'' of the network &
    Medium & High & Default-deny at internal boundaries; documented
    exceptions with expiration; ETH isolation zones \\
\addlinespace
Teaching labs (offensive security courses) impacting production &
    Low & High & Isolated cyber-ranges; rapid reset; strict
    egress policies \\
\addlinespace
Insufficient \gls{gdpr} compliance for AI access outside the EU &
    Medium & High & Researcher-centred accountability per
    Section~\ref{sec:safeguards-ai}: PI data classification +
    IT-service register + cumulative \emph{Schrems~II} regime for
    personal-data transfers + EU-first preference + periodic DPO
    audit~\cite{schrems2020} \\
\bottomrule
\end{tabularx}
\end{table}

\subsection{Monitoring KPIs}

\begin{table}[H]
\centering
\caption{Key performance indicators}
\label{tab:kpi}
\small
\begin{tabularx}{\textwidth}{X X}
\toprule
\textbf{KPI} & \textbf{Measure} \\
\midrule
Reduction of admin enclave exposure & Services reachable
    from non-admin zones \\
\addlinespace
Teaching lab rebuild time & Reimage/reset cycle;
    incidents resolved by reset \\
\addlinespace
Research enclave onboarding rate & Chartered enclaves;
    request$\rightarrow$approval delay \\
\addlinespace
Lateral movement containment & Post-mortems: did the compromise
    cross zone boundaries? \\
\addlinespace
DTN throughput and diagnostic time & Transfer performance, perfSONAR
    deviations \\
\addlinespace
Art. 32 evidence production & Planned and completed control tests,
    audits, exercises \\
\bottomrule
\end{tabularx}
\end{table}

% ============================================================================
\section{Propositions}
\label{sec:propositions}
% ============================================================================

\textbf{Implementation strategy and cost control.} The propositions
below are ordered so that \emph{governance and quick wins} come first
(risk-based approach, federated committee, charters): these require
little or no additional infrastructure and deliver immediate clarity
of responsibilities. \emph{Pilot deployment} (one department, one zone)
then validates the model at contained cost before any institution-wide
rollout. The heavy use of shared services (the \gls{nren}, its continental
research backbone --- in Europe, \gls{geant} ---, REN-ISAC,
open-source tooling) and the absence of a single ``critical path''
vendor reduce both implementation risk and total cost of ownership.
The cost-benefit and feasibility discussion in
Section~\ref{sec:couts} and the phased transition in
Section~\ref{sec:transition} should be read as an integral part of
these propositions.

\bigskip
\begin{enumerate}
    \item \textbf{Adopt a risk-based approach} (NIST CSF 2.0,
    \gls{gdpr} Art.~32) rather than an approach mainly driven by uniform restrictions.
    Each zone receives controls proportionate to its risk level.

    \item \textbf{Segment the network into 7 zones}: administrative
    (high security), teaching, research (flexible), \gls{sciencedmz},
    personal \gls{byod}, IoT, sensitive data enclave. These zones
    constitute the operational architecture through which the
    institution implements the proportionate measures required by the
    national NIS2 transposition; per-zone compliance evidence
    aggregates into the institutional annual maturity report
    submitted to the national supervisory authority.

    \item \textbf{Create a personal \gls{byod} zone} with Internet access
    and progressive interoperability with other zones depending on
    authentication.

    \item \textbf{Guarantee access to AI cloud services} from the
    research zone under a \emph{researcher-centred accountability
    model}. The Principal Investigator classifies data and selects
    providers proportionately; personal-data transfers outside the
    EEA follow the cumulative \emph{Schrems~II} regime (\gls{scc} +
    supplementary measures)~\cite{schrems2020}. International
    collaboration with non-EU providers is permitted on a
    \emph{cooperative declaration} basis: the central IT service
    registers the access on PI request, the register is available
    for periodic \gls{dpo} audit; EU-first is a routing preference,
    not a restriction (operational detail in
    Section~\ref{sec:safeguards-ai}; advocacy positioning in
    Section~\ref{sec:advocacy}, Position~D.5).

    \item \textbf{Establish a Digital Policy Commission} chaired by
    an elected academic, with representatives from departments and
    interdisciplinary centres, responsibility charters, and a service
    catalogue.

    \item \textbf{Adopt open-source and open-science tools as default
    institutional choices}, reducing vendor lock-in, increasing
    auditability and transparency, and aligning with public-sector
    open-source strategies (e.g., the European Commission's
    open-source strategy in the EU).

    \item \textbf{Formalize autonomy} through enclave charters
    (ETH Zurich / Oxford model) and Data Stewards in
    laboratories (EPFL model).

    \item \textbf{Grant administrative autonomy to research-only units,
    research-and-teaching departments, and their laboratories} within
    their dedicated \gls{vlan}s, with minimum security standards defined
    centrally but freedom of local implementation. Administrative
    autonomy is exercised within the institutional minimum standards
    required by the national NIS2 transposition (mirroring NIS2
    Article~21(2)) and opposable to
    all zones; enclave charters formalise these standards and their
    verification by internal audit feeds the maturity report
    submitted to the national supervisory authority. \textbf{Autonomy
    is contextual, not exempt.}

    \item \textbf{Deploy a \gls{sciencedmz}} for high-performance research
    data transfers, following the ESnet pattern adopted at over 100
    university campuses with NSF \emph{Campus Cyberinfrastructure} (CC*) funding~\cite{esnet-history,nsf-cc}.

    \item \textbf{Implement federated identity} via
    \gls{eduroam} and \loc{identity-federation-profile}, enabling inter-institutional
    \gls{sso} and researcher mobility --- where a national identity
    federation operates; otherwise deploy the standards-based
    identity layer with interoperable federation metadata and defer
    participation until one exists (Technical Reference, Ch.~6).

    \item \textbf{Offer secure alternatives} rather than
    banning: OpenStack private cloud, \gls{vdi} for teaching,
    Kubernetes for research environments, rather than
    blocking cloud access or software installation.

    \item \textbf{Participate in security communities} (\gls{first}
    globally, \gls{geant} Security in Europe, and the equivalent regional
    body --- e.g.\ REN-ISAC in the US, JRES in France) to share threat
    intelligence and best practices.

    \item \textbf{Implement \gls{zt} progressively}:
    start with identity (universal \gls{mfa}), then device
    posture, then micro-segmentation, adapting requirements
    to the academic context.

    \item \textbf{Demonstrate \gls{nis2} essential-entity compliance} by
    adopting an architecture producing per-zone evidence of risk
    management, incident containment and Article~21 proportionate
    measures, supporting the institution's annual compliance and
    maturity report to the national supervisory authority.
    Per-zone evidence consolidates into an institution-level maturity
    dashboard, based on an ISO/IEC~27005-aligned risk-analysis method
    (such as a nationally provided open-source tool),
    producing the single annual score submitted to the supervisory
    authority pursuant to the national NIS2 transposition.

    \item \textbf{Draw inspiration from reference universities}:
    Oxford, ETH Zurich, Michigan, Rutgers, Cambridge, MIT, Stanford,
    CERN, Penn State, NC State, and Wisconsin-Madison demonstrate that
    security and academic freedom are compatible.
\end{enumerate}

\medskip
\noindent The propositions above, taken together, constitute a
\keyword{continuous-improvement mechanism}: each proposition
contributes to the per-zone evidence base and to the federated
governance dispositive whose institutional artefact is the annual
compliance and maturity report to the
national supervisory authority pursuant to the national NIS2
transposition. The continuous-improvement
framing is aligned with cybersecurity governance best practice
(NIST CSF~2.0 GV function~\cite{nist-csf}, ISO/IEC~27001:2022
clause~10) and with the maturity-cycle expectations of NIS2
essential entities.

% ============================================================================
\section{Independent Commission for Reasonable Adjustments to Digital Policies}
\label{sec:commission-ajustements}
% ============================================================================

\epigraph{``Where, after all, do universal human rights begin? In small places, close to home~[\ldots] Unless these rights have meaning there, they have little meaning anywhere.''}{--- Eleanor Roosevelt, \textit{In Your Hands}, 1958~\cite{roosevelt1958}}

Independent of the IT security policy adopted by the institution,
a mechanism allowing any employee to obtain
\keyword{reasonable adjustments} for personal, medical or
professional reasons is a structural requirement. This mechanism
is to be guaranteed by an \textbf{independent commission}
operating outside the operational chain of command of the IT
services.

\subsection{Legal and Normative Foundations}

\subsubsection{European Equal Treatment and Disability Law}

The concept of \emph{reasonable adjustment} (\emph{reasonable
accommodation}) is a pillar of European anti-discrimination law:

\begin{itemize}
    \item \textbf{Directive 2000/78/EC}~\cite{directive-2000-78} establishing
    a general framework for equal treatment in
    matters of employment and occupation requires employers to take
    \emph{``appropriate measures''} to enable a person with a disability
    to access, exercise, and advance in employment (Article~5). The
    CJEU has interpreted `disability' to include some long-term
    limitations, including certain chronic or long-lasting illnesses,
    but not sickness as such and not purely temporary
    conditions~\cite{cjue-chacon-navas, cjue-ring-skouboe}.

    \item The \textbf{United Nations Convention on the Rights of
    Persons with Disabilities} (CRPD)~\cite{un-crpd}, ratified by the EU and
    its member states (and, more broadly, by the relevant jurisdiction),
    defines ``discrimination on the basis of disability''
    to include the \emph{denial} of reasonable accommodation (Article~2,
    Definitions) and establishes the substantive duty of equality and
    non-discrimination, including reasonable accommodation, in
    Article~5.

    \item \textbf{\loc{national-equal-treatment-law}}\loc{cite-equal-treatment-law}
    integrates the obligation of
    reasonable adjustments into national law.
\end{itemize}

\subsubsection{Extension Beyond Disability: the Concept of Reasonable Adjustment of Digital Policies}

\textbf{Note:} the extension of the reasonable adjustment concept
beyond disability is proposed here as \emph{institutional governance
best practice} inspired by the legal model, not as a legal obligation
under Directive 2000/78/EC or the CRPD --- which specifically concern
disability. The governance rationale is that frameworks such as
COBIT~2019 (EDM05), NIST CSF~2.0 (GV function) and ISO~27001
(Clause~4.2) all require integrating stakeholder needs into IT
governance, providing an independent institutional basis for this
mechanism.

\medskip

While the original legal framework primarily concerns disability,
\textbf{good management practice} and \textbf{IT governance
frameworks} recognise that uniform digital policies
can create disproportionate obstacles for various
reasons:

\begin{description}
    \item[Medical reasons:] an employee suffering from musculoskeletal
    or visual disorders may need specific equipment
    or software not included in the standard catalogue;
    a researcher undergoing treatment requiring adapted telework may
    need a particular \gls{vpn} access.

    \item[Personal reasons:] a single parent providing
    support for a child with a disability may need
    remote access flexibility incompatible with a restrictive
    \gls{byod} policy; an employee acting as a family caregiver may
    need specific communication tools.

    \item[Professional reasons:] a cybersecurity researcher
    needs access to malware analysis environments; an
    artificial intelligence instructor must access non-European cloud APIs;
    a data science researcher requires
    local administrative privileges (root/admin access) on their workstation.
\end{description}

The \textbf{NIST Cybersecurity Framework 2.0}~\cite{nist-csf}, in
its \emph{Govern} (GV) function, insists on integrating
\emph{``stakeholder needs''} into the cybersecurity
strategy. The \textbf{COBIT 2019}~\cite{cobit} framework emphasizes
that IT governance must serve \emph{``stakeholder
needs''} and not constrain them without
justification. \textbf{ISO 27001:2022}~\cite{iso27001}, in its
Clause~4.2, explicitly requires understanding the
\emph{``needs and expectations of interested parties''} as a prerequisite
for the \gls{isms}.

\subsubsection{Infringement of Fundamental Freedoms and Digital Conscientious Objection}

Beyond medical, personal, and professional reasons, a
fourth category of grounds for adjustment must be recognised: the
situation of an employee whose \textbf{fundamental freedoms} are
infringed disproportionately by a digital policy, or
who feels, in conscience, that this policy constitutes a violation
of rights they consider inalienable. This ground, which we
call \keyword{digital conscientious objection}, is based on
a converging body of European and international fundamental rights.
No court has yet applied Article~9 to digital workplace policies;
however, the \emph{Bayatyan} criteria --- a serious and insurmountable
conflict between an obligation and a sincerely held conviction ---
provide an arguable framework, particularly where restrictions compel
actions that conflict with deeply held professional or ethical
convictions.

\paragraph{The fundamental freedoms at stake.}

Restrictive university digital policies simultaneously engage
several fundamental rights. In EU-bound jurisdictions these are
grounded in the \textbf{Charter of Fundamental Rights of the
EU}~\cite{eu-charter} and the \textbf{European Convention on Human
Rights}~\cite{echr}; in other jurisdictions the adopter's equivalent
constitutional or human-rights instruments apply, with the
international consensus (below) as the jurisdiction-neutral anchor:

\begin{description}[style=unboxed]
    \item[Academic freedom (EU Charter, Art.~13):] \emph{``The arts
    and scientific research shall be free of constraint. Academic freedom shall be
    respected.''} This right is self-standing within
    the Charter. The \gls{ecthr} linked it to Article~10 in the judgment
    \emph{Sorgu\c{c} v.\ Turkey}~\cite{sorguc2009}, which recognises the
    freedom of academics to \emph{``freely express their opinion
    about the institution or system in which they work and to distribute
    knowledge and truth without restriction''}. The judgment
    \emph{Mustafa Erdo\u{g}an v.\ Turkey}~\cite{erdogan2014},
    in which a concurring opinion (Judges Saj\'o, Vu\v{c}ini\'c
    and Kuri\v{s}) proposes that \emph{``academic discourse should
    enjoy the highest level of protection''}.

    \item[Freedom of expression and information (EU Charter, Art.~11;
    ECHR, Art.~10):] includes \emph{``the freedom to receive or
    impart information or ideas without interference by
    public authorities and regardless of frontiers''}. The \gls{ecthr}
    characterised the Internet as \emph{``one of the principal means by
    which individuals exercise their right to freedom to receive
    and impart information''} (\emph{Yildirim v.\
    Turkey}~\cite{yildirim2012}). Systematic filtering and blocking
    were held incompatible with Article~11 of the
    Charter by the \gls{cjeu} in \emph{Scarlet Extended v.\
    SABAM}~\cite{scarlet2011}.

    \item[Freedom of thought and conscience (EU Charter, Art.~10;
    ECHR, Art.~9):] the Grand Chamber of the \gls{ecthr} recognised in
    the judgment \emph{Bayatyan v.\ Armenia}~\cite{bayatyan2011} that
    conscientious objection is an autonomous right protected by
    Article~9. While the case concerned military service, the
    \emph{Bayatyan} criteria --- a serious and insurmountable conflict
    between an obligation and a sincerely held conviction --- may
    provide an arguable framework for other contexts. The judgment
    \emph{Eweida v.\ United Kingdom}~\cite{eweida2013} extended the
    protection of Article~9 in the professional context.

    \item[Respect for private life and communications (EU Charter,
    Art.~7 and 8; ECHR, Art.~8):] the Grand Chamber judgment
    \emph{B\u{a}rbulescu v.\ Romania}~\cite{barbulescu2017} established a
    \textbf{six-criteria proportionality test} for
    digital surveillance in the workplace and laid down the principle that
    \emph{``restrictions cannot reduce to zero private social
    life in the workplace''}.

    \item[Dignity at work (EU Charter, Art.~1 and 31):] Article~31
    guarantees every worker \emph{``working conditions which
    respect their health, safety, and \textbf{dignity}''}. Dignity
    is explicitly listed as a component of fair
    working conditions.
\end{description}

\paragraph{The international consensus on digital rights.}

\begin{itemize}[nosep]
    \item Resolution A/HRC/RES/20/8 of the \textbf{UN Human Rights
    Council}~\cite{unhrc2012} (2012) affirms for the
    first time that \emph{``the same rights that people have
    offline must also be protected online''}.
    \item Resolution A/HRC/RES/32/13~\cite{unhrc2016} (2016)
    \textbf{condemns} measures aimed at intentionally preventing or disrupting
    access to information online.
    \item Resolution A/HRC/RES/47/16~\cite{unhrc2021} (2021),
    adopted without a vote (by consensus) and co-sponsored by some
    70 states, condemns Internet shutdowns and
    online censorship.
    \item The \textbf{UNESCO Recommendation on Science and
    Scientific Researchers}~\cite{unesco2017} (2017, 195 countries)
    affirms that researchers must \emph{``work in a
    spirit of intellectual freedom''} and have
    \emph{``free access to libraries and other sources
    of information''}.
    \item The \textbf{Magna Charta Universitatum}~\cite{magna-charta}
    (Bologna, 1988, signed by university rectors from around the world, revised in 2020)
    states that \emph{``research and teaching must be
    intellectually and morally independent of all political
    authority and economic power''}.
\end{itemize}

\paragraph{The proportionality requirement.}

Article~52(1) of the EU Charter requires that \textbf{any
limitation} of a fundamental right:
\begin{enumerate}[nosep]
    \item be \emph{provided for by law};
    \item \emph{respect the essential content} of the right;
    \item be \emph{proportionate}, necessary, and \emph{genuinely}
    meet an objective of general interest.
\end{enumerate}

The \gls{cjeu} has invalidated disproportionate digital restrictions in
three major judgments: \emph{Digital Rights Ireland}~\cite{dri2014}
(invalidation of the data retention directive),
\emph{Schrems I}~\cite{schrems2015} and
\emph{Schrems II}~\cite{schrems2020} (invalidation of
EU-US transfer agreements for disproportionate surveillance). These judgments
establish that an \textbf{indiscriminate and
disproportionate digital restriction is illegal} under EU law, even when it
pursues a legitimate security objective.

\paragraph{Technostress and moral injury as documented
harms.}

The scientific literature documents two forms of harm
caused by restrictive digital policies:

\begin{itemize}
    \item \keyword{Technostress}~\cite{brod1984, tarafdar2007,
    tarafdar2015, salanova2013, latorre2019} is an empirically
    established occupational risk (systematic review of 105~studies),
    characterised by anxiety, fatigue, and loss of effectiveness.
    The absence of autonomy in technology use is a positive predictor of
    technostress~\cite{salanova2013} --- a finding that, while not
    directly studied in the specific context of restrictive IT policies,
    provides a plausible mechanism linking rigid digital environments to
    occupational stress. Technostress reduces job satisfaction and
    organisational commitment~\cite{ragunathan2008}.

    \item \keyword{Moral injury} (\emph{moral injury}), defined
    as the trauma resulting from \emph{``acts that transgress
    deeply held moral beliefs and
    expectations''}~\cite{litz2009}, was systematically reviewed
    in the \emph{British Journal of
    Psychiatry}~\cite{williamson2018}, which called for extension
    of the concept to civilian occupational contexts. While no study has directly
    measured moral injury resulting from IT policies, the structural
    conditions identified in the occupational literature --- compulsion to
    act against professional values, perceived futility of action, and
    absence of institutional recourse --- can be present when restrictive
    digital policies force academics to choose between professional
    obligations and policy compliance. A qualitative study
    documented the moral injury of university
    staff confronted with institutional policies
    experienced as a betrayal of their professional
    values during the COVID-19 crisis~\cite{hanna2023}. The concept of
    \emph{``chronic moral injury''}~\cite{abadal2022}
    applies to situations where a professional aspiring to
    excellence must operate within an institution whose
    practices hinder the exercise of their competencies.

\smallskip\noindent
\textbf{Note on proportionality:} the severity of moral stakes in IT
policy contexts is clearly lower than in military or pandemic contexts
where the original concept was developed. The argument here is not one
of equivalence but of \emph{structural analogy}: the mechanism
(institutional compulsion against professional values, absence of
recourse) may operate in attenuated form when digital policies
systematically prevent academics from fulfilling their professional
obligations.
\end{itemize}

\paragraph{Summary: the eligibility of requests based on
fundamental freedoms.}

The legal argument is structured as follows:

\begin{enumerate}
    \item Restrictive university digital policies
    \textbf{engage} simultaneously academic freedom (Art.~13),
    freedom of expression and information (Art.~11), freedom of
    conscience (Art.~10), the right to private life and
    communications (Art.~7-8) and dignity at work (Art.~1, 31)
    of the EU Charter;
    \item Any restriction of these rights must pass the
    \textbf{proportionality test} of Article~52(1), which
    implies a case-by-case assessment --- not an automatic refusal;
    \item \textbf{Indiscriminate and
    disproportionate} restrictions are illegal (\emph{Digital Rights Ireland},
    \emph{Scarlet Extended});
    \item Disproportionate restrictions can cause
    \textbf{documented harm} (technostress, moral injury);
    \item An employee who experiences that a digital policy
    \emph{disproportionately infringes on their fundamental freedoms}
    --- whether through deep conviction (conscientious objection,
    Art.~9/10), through infringement of their academic freedom
    (Art.~13), or through harm to their dignity and health (Art.~31) ---
    has a \textbf{defensible legal basis} to request
    a reasonable adjustment;
    \item For requests grounded in disability or other protected
    characteristics (Directive 2000/78/EC), the burden shifts to the
    institution to demonstrate that the refusal is proportionate. For
    other grounds (personal, professional), the requester must provide
    a documented justification, after which the institution must
    demonstrate that refusal is proportionate to an identified security
    risk --- not merely administratively inconvenient.
\end{enumerate}

\textbf{Important:} this ground does not open an unlimited right to
bypass any security rule. It requires an \emph{individualized
proportionality assessment} by an independent body,
in accordance with the consistent case law of the \gls{cjeu} and the \gls{ecthr}.
The adjustment granted may be an alternative technical solution
(different network zone, dedicated enclave, specific \gls{vpn}) that reconciles
the individual's need with the legitimate security objectives.

\subsection{Why an \emph{independent} commission?}

\subsubsection{The problem of internal conflict of interest}

When adjustment requests are handled by the IT department
itself --- the one that designs and enforces the policy --- a
\textbf{structural conflict of interest} arises:

\begin{itemize}[nosep]
    \item The IT department has an operational interest in minimizing
    exceptions, which increase its workload and support
    surface;
    \item The requester is in a position of technical subordination
    vis-\`a-vis the department that evaluates their request;
    \item The absence of external oversight favors arbitrary decisions
    or unmotivated refusals;
    \item The literature on organisational governance shows that
    self-assessment mechanisms without external oversight tend toward
    \emph{``status quo bias''}~\cite{samuelson-zeckhauser}.
\end{itemize}

As De~Haes and Van~Grembergen~\cite{dehaes2015} emphasize, effective
IT governance relies on \emph{``structures, processes
and relational mechanisms''} that separate strategic decision-making from
operational execution. Entrusting the arbitration of adjustments to
the department that enforces the policy violates this fundamental principle.

\subsubsection{The university Ombudsman model}

Many universities have an \textbf{Ombudsman}
(\emph{independent mediator}) tasked with handling conflicts between
individuals and the institution in an impartial manner. The International
Ombudsman Association (IOA)~\cite{ioa-standards} defines four
fundamental standards: \textbf{independence}, \textbf{neutrality},
\textbf{confidentiality} and \textbf{informality}. These principles
apply directly to a digital policy adjustment commission.

\subsection{The precedent of an existing reasonable-adjustments committee}

\loc{existing-adjustments-committee-precedent} that supports
students and staff with disabilities, learning difficulties,
chronic illnesses, or other specific
needs. Such a committee typically:

\begin{itemize}[nosep]
    \item Reviews adjustment requests and recommends practical
    adjustments in collaboration with the relevant institutional actors;
    \item Covers a wide range of measures: additional time for
    examinations, alternative submission formats, adaptation of
    teaching materials, individualised assessment modalities;
    \item Operates on the principle that the \textbf{general framework
    should adapt to the individual} when a legitimate need is
    documented.
\end{itemize}

\textbf{If the institution already operates a mechanism for adjusting
academic practices to individual needs --- including the design and
assessment of courses --- it is coherent to extend an analogous
principle to digital policies}, which are a tool in service of
the university's missions.

\subsection{Proposal: Independent Commission for Reasonable Adjustments of Digital Policies (ICRADP)}

\subsubsection{Mission}

The \gls{icradp}'s mission is to ensure that \textbf{every employee}
--- whether in a research-only unit, a research-and-teaching department,
a teaching/training component, or an administrative service ---
can obtain a reasonable adjustment to any university digital policy
when that policy creates a disproportionate obstacle
to the exercise of their functions or infringes on
their fundamental rights, for reasons of:
\begin{itemize}[nosep]
    \item \textbf{Medical}: disability, chronic illness, temporary
    health condition;
    \item \textbf{Personal}: family situation, caregiver
    obligations, documented individual circumstances;
    \item \textbf{Professional}: specific needs related to
    research, teaching, or the exercise of a particular
    function;
    \item \textbf{Infringement of fundamental freedoms}: disproportionate
    infringement of academic freedom, freedom of expression
    and information, freedom of conscience, dignity at
    work, or documented harm (technostress, moral injury)
    caused by the digital policy
    (see Section~\ref{sec:commission-ajustements},
    \S\emph{Infringement of fundamental freedoms}).
\end{itemize}

\subsubsection{Fundamental principles}

\begin{enumerate}
    \item \textbf{Independence:} the commission reports to \loc{executive-leadership}
    and not to the IT department or any operational service. Its members
    may not receive instructions from any department regarding the
    handling of requests.

    \item \textbf{Pluralist composition:}
    \begin{itemize}[nosep]
        \item An independent chair (tenured teaching-researcher
        without a management function);
        \item A representative of the teaching staff elected by peers;
        \item A representative of administrative and technical staff;
        \item The Data Protection Officer (\gls{dpo}), as a
        legal expert;
        \item An external expert in digital ergonomics or labour
        law (advisory member);
        \item The CISO is \emph{heard} but is \textbf{not}
        a voting member, to preserve independence.
    \end{itemize}

    \item \textbf{Universal right of referral:} any employee may refer
    a matter to the commission directly, without going through their
    hierarchy or the IT department.

    \item \textbf{Processing deadline:} the commission rules within a
    maximum period of \loc{icradp-decision-deadline}. In case of a documented medical or
    professional emergency, a provisional adjustment may be
    granted within \loc{icradp-emergency-deadline}.

    \item \textbf{Reasoned decisions:} any refusal of adjustment
    must be justified in writing, demonstrating that the requested
    adjustment would constitute a \emph{``disproportionate burden''} within
    the meaning of Directive 2000/78/EC or would create a
    \emph{documented and quantified} security risk (and not merely invoked).

    \item \textbf{Right of appeal:} the requester may appeal
    to \loc{executive-leadership} in case of refusal.

    \item \textbf{Transparency:} the commission publishes an
    anonymized annual report indicating the number of requests, the types
    of adjustments granted and refused, and the processing times.

    \item \textbf{Security safeguards:}
    \begin{itemize}[nosep]
        \item A security impact assessment is required before any binding
        decision --- the CISO provides a written analysis of the risks
        associated with the requested adjustment;
        \item Certain categories of adjustment are excluded \emph{a priori}:
        exemptions from \gls{mfa}, access to the
        sensitive data enclave without standard controls, and any
        adjustment that would create an unmonitored network path to
        the infrastructure core;
        \item The commission must include at least one member with
        demonstrated competence in information security or digital risk.
    \end{itemize}
\end{enumerate}

\subsubsection{Examples of reasonable adjustments}

\begin{table}[H]
\centering
\caption{Examples of reasonable adjustments to digital policies}
\label{tab:ajustements}
\small
\begin{tabularx}{\textwidth}{>{\raggedright\arraybackslash}X >{\raggedright\arraybackslash}X >{\raggedright\arraybackslash}X}
\toprule
\textbf{Reason} & \textbf{Obstacle} & \textbf{Possible adjustment} \\
\midrule
Medical (visual impairment) & Unauthorized accessibility software &
    Installation authorization with IT support \\
\addlinespace
Medical (RSI) & Standard keyboard/mouse imposed &
    Authorization of ergonomic BYOD equipment \\
\addlinespace
Personal (caregiver) & Restricted remote access &
    Extended VPN with enhanced \gls{mfa} \\
\addlinespace
Professional (AI research) & Cloud AI APIs blocked &
    Access from dedicated research zone \\
\addlinespace
Professional (security teaching) & Pentest tools prohibited &
    Isolated cyber-range with charter \\
\addlinespace
Professional (root admin) & Admin rights denied &
    Granting of admin rights (ETH model) \\
\addlinespace
Academic freedom (Art.~13) & Web filtering blocking scientific
    sources & Unfiltered access from research zone \\
\addlinespace
Freedom of information (Art.~11) & Blocking of major cloud-based AI
    services & Access via research zone with traceability \\
\addlinespace
Conscience / dignity (EU Charter Art.~10, 31) & Monitoring perceived as
    an affront to dignity & Reduction of monitoring to the
    strict minimum, BYOD zone \\
\addlinespace
Documented technostress & Restrictions causing anxiety and loss
    of efficiency & Individualized easing with follow-up \\
\bottomrule
\end{tabularx}
\end{table}

\subsection{Justification from the literature and international best practices}

Several converging arguments support the creation of such a
commission:

\begin{enumerate}
    \item \textbf{International law requires it for disability:}
    the CRPD~\cite{un-crpd} and Directive 2000/78/EC~\cite{directive-2000-78}
    mandate reasonable adjustments. Yet digital policies
    \emph{are} working conditions --- their rigidity may constitute
    indirect discrimination within the meaning of the directive.

    \item \textbf{The risk-based approach recommends it:}
    if security must be \emph{proportionate to risk}
    (\gls{gdpr} Art.~32~\cite{gdpr}), then the refusal of adjustment must be
    proportionate to the actual risk --- which requires an independent
    assessment, not an automatic refusal by the IT department.

    \item \textbf{The best universities practice it:}
    \begin{itemize}[nosep]
        \item Oxford, MIT and Stanford have governance structures in which
        academic representatives participate in IT policy decisions,
        providing institutional channels for addressing disproportionate
        restrictions;
        \item ETH Zurich~\cite{eth-baseline} has institutionalized
        exception pathways (self-assessment, registration, temporary
        exemptions, isolation zones);
        \item The Harvard Faculty of Arts and Sciences has an
        \emph{IT Governance Committee} of academics that arbitrates
        conflicts between IT policy and academic
        needs~\cite{harvard-it-governance}.
    \end{itemize}

    \item \textbf{Organizational psychology confirms it:}
    Deci and Ryan~\cite{deci-ryan-sdt} show that perceived autonomy
    is a fundamental psychological need. Rigid digital policies
    without recourse diminish the sense of autonomy and
    competence, reducing intrinsic motivation and increasing
    turnover --- a major risk for a university that depends on
    the attractiveness of its talent.

    \item \textbf{IT governance prescribes it:}
    COBIT 2019~\cite{cobit} identifies the process
    \emph{``Ensure Stakeholder Engagement''} (EDM05) as a governance
    objective. The absence of a recourse mechanism for
    employees violates this principle.
\end{enumerate}

\subsection{Integration with the proposed architecture}

The \gls{icradp} is not in contradiction with the segmented architecture
proposed in this dossier --- it is its \textbf{indispensable human
complement}. The multi-zone architecture structurally reduces
the need for individual adjustments (most needs are
covered by the appropriate zone), but it cannot cover all
cases. The \gls{icradp} ensures that:

\begin{itemize}[nosep]
    \item Individual cases not covered by the zone division
    are handled fairly and promptly;
    \item The IT department is not overloaded by ad hoc exception
    requests --- the commission centralizes and formalizes the process;
    \item The university demonstrates its compliance with labour law and
    anti-discrimination law;
    \item The digital policy remains a \emph{tool in service of
    people}, and not a rigid system imposed on individuals
    without recourse.
\end{itemize}

% ============================================================================
\section{Conclusion}
\label{sec:conclusion}
% ============================================================================

\epigraph{``One resists the invasion of armies; one does not resist the invasion of ideas.''}{--- Victor Hugo, \textit{Histoire d'un Crime}~\cite{hugo1877}}

The international state of the art shows a clear trend: the most
high-performing and most secure universities in the world generally do \emph{not}
rely primarily on a uniform restriction policy. They tend to adopt
\keyword{segmented and federated architectures} that:

\begin{itemize}[nosep]
    \item Strongly protect sensitive data in dedicated zones;
    \item Offer great flexibility to researchers and instructors
    in adapted zones;
    \item Use monitoring and detection rather than
    prohibitions as the first line of defence on research
    networks;
    \item Rely on identity federation and mutual
    security communities.
\end{itemize}

The fundamental principle of \keyword{dependability} ---
segmentation combined with Trusted Activity Contexts --- confirms this
approach: it is by multiplying barriers, zones, and validated activity
environments, and not by reducing everything to a single perimeter,
that the best resilience is achieved.

In the context of 2026, three additional dimensions make this
change even more urgent:
\begin{itemize}[nosep]
    \item \textbf{Artificial intelligence} increasingly requires access
    to global cloud services to remain competitive in research and
    teaching;
    \item \textbf{Staff \gls{byod}} is a widespread reality
    that benefits from being integrated into the architecture rather than resisted;
    \item The \textbf{European regulatory framework} (\gls{gdpr}, \gls{nis2}, AI Act)
    requires measures \emph{proportionate to risk} and a capacity for
    \emph{evidence} --- exactly what a multi-zone architecture
    with federated governance produces.
\end{itemize}

A security policy that significantly hinders researchers from researching and
instructors from teaching risks weakening the university rather than protecting it.
Security should serve as an \emph{enabler} (\emph{enabling security}),
not as an obstacle. Proven solutions exist and are already
deployed in many leading universities worldwide. Adopting them would
strengthen both security and the university's core missions.

\medskip
\noindent
The \textbf{feasibility and cost} of the proposed measures have been
addressed explicitly: implementation is \emph{phased} (governance
$\rightarrow$ pilot $\rightarrow$ rollout), relies on existing
\gls{nren} and continental research-backbone ecosystems (in Europe,
\gls{geant}) and open-source or shared tools where possible,
and concentrates investment on the administrative enclave rather than
on the whole perimeter. This makes the approach both \emph{implementable}
within typical university budgets and \emph{proportionate} to risk, as
required by \gls{gdpr} Article~32 and good governance practice.

\medskip
\noindent
Ultimately, a university requires administration, but is not a corporation;
university governance must optimise for freedom for ideas and persons,
physical or digital. The digital infrastructure should reflect this
fundamental distinction.
